% Alternate review version: front matter + intro + contributions + technical content only.
% Keeps the original title/author/abstract/intro unchanged (via \input) and preserves theorem/proof statements.

\documentclass[11pt,a4paper]{article}
\newif\ifjeicblind
\jeicblindfalse

\usepackage[utf8]{inputenc}
\usepackage[T1]{fontenc}
\usepackage{graphicx}
\usepackage{amsmath, amssymb, amsthm}
\usepackage{booktabs}
\usepackage{algorithm}
\usepackage{algorithmic}
\usepackage{enumitem}
\usepackage{url}
\usepackage[round,authoryear]{natbib}
\usepackage[hmargin=2.5cm,vmargin=3cm]{geometry}
\usepackage{hyperref}

\graphicspath{{../paper_artifacts/figures/}{../}}

% JEIC double-blind: clear author/keywords in PDF metadata; wrap acknowledgments and data-availability in \jeicAck{...} and \jeicDataAvailability{...}
\ifjeicblind
  \hypersetup{pdfauthor={},pdfkeywords={}}
\fi

% Springer-like keywords (simple paragraph)
\newcommand{\keywords}[1]{\par\noindent\textbf{Keywords:} #1}

% Theorem-like environments
\theoremstyle{plain}
\newtheorem{theorem}{Theorem}[section]
\newtheorem{lemma}[theorem]{Lemma}
\newtheorem{proposition}[theorem]{Proposition}
\newtheorem{corollary}[theorem]{Corollary}

\theoremstyle{definition}
\newtheorem{definition}[theorem]{Definition}
\newtheorem{assumption}[theorem]{Assumption}

\theoremstyle{remark}
\newtheorem{remark}[theorem]{Remark}
\newtheorem{example}[theorem]{Example}

% Compact notation
\newcommand{\R}{\mathbb{R}}
\newcommand{\N}{\mathbb{N}}
\newcommand{\E}{\mathbb{E}}
\newcommand{\Prob}{\mathbb{P}}
\newcommand{\Var}{\text{Var}}
\newcommand{\Cov}{\text{Cov}}
\newcommand{\Law}{\mathcal{L}}
\newcommand{\ind}{\mathbf{1}}
\newcommand{\norm}[1]{\left\|#1\right\|}
\newcommand{\abs}[1]{\left|#1\right|}
\newcommand{\inner}[2]{\left\langle#1,#2\right\rangle}
\newcommand{\argmin}{\operatorname{argmin}}
\newcommand{\argmax}{\operatorname{argmax}}
\newcommand{\diag}{\operatorname{diag}}
\newcommand{\trace}{\operatorname{tr}}
\newcommand{\rank}{\operatorname{rank}}
\newcommand{\sign}{\operatorname{sign}}
\newcommand{\supp}{\operatorname{supp}}
\newcommand{\sigeta}{\sigma_\eta}

\newcommand{\MFG}{\text{MFG}}
\newcommand{\HJB}{\text{HJB}}
\newcommand{\FPK}{\text{FPK}}
\newcommand{\CARA}{\text{CARA}}
\newcommand{\LQG}{\text{LQG}}
\newcommand{\REE}{\text{REE}}

% Acknowledgments: wrap in \jeicAck{...} so omitted in blind mode.
\newcommand{\jeicAck}[1]{\ifjeicblind\else #1\fi}
% Data/code availability in front matter: use \jeicDataAvailability{...}; in blind mode it is replaced by the standard phrase.
\newcommand{\jeicDataAvailability}[1]{\ifjeicblind Replication package will be provided upon acceptance.\else #1\fi}

\title{Mean Field Game Modeling of Overconfidence in Financial Markets}
\ifjeicblind
  \author{Anonymous Author(s)}
\else
  \author{Lucas~Sneller\\[0.5em]\normalsize Miami University, Oxford, OH 45056, USA}
\fi
\date{}

\begin{document}
\maketitle
\thispagestyle{empty}

\begin{abstract}
This paper develops a mean field model of overconfidence in a market with flow-driven price impact and inventory carryover. A continuum of investors filters an unobserved fundamental from noisy private signals and trades against an endogenous price that mean-reverts toward fundamentals and responds to aggregate order flow; overconfidence enters as misperceived private-signal precision, generating heterogeneous filtered beliefs. Our main equilibrium object is a discrete-time \emph{myopic mean-field equilibrium}: given a conjectured mean trading rate, each agent chooses next-step inventory to maximize one-step \CARA utility under conditionally Gaussian price increments. This yields a linear feedback rule and reduces mean-field consistency to a scalar stock--flow fixed point at each time step. We provide an explicit per-step contraction condition guaranteeing existence, uniqueness, and stability of the resulting discrete-time dynamics and particle approximation. In simulations across bull/bear regimes and a targeted joint grid over anchoring, impact, and noise trading, microstructure parameters govern mispricing levels, whereas overconfidence primarily increases disagreement and trading intensity with limited effect on price-level mispricing in the baseline anchored-impact regime. A finite-horizon continuous-time \CARA/\LQG\ solution is reported as a benchmark, and an identification/moment-matching vignette connects the model's parameters to price and volume moments.
\end{abstract}

\section{Introduction}
\label{sec:introduction}
% Section 1: Introduction

Overconfidence---the tendency to overestimate the precision of one's information---is a well-documented behavioral bias in financial markets and is empirically associated with excessive trading and episodes of mispricing \cite{Odean1998,BarberOdean2001}. A central modeling challenge is to combine heterogeneous (and potentially biased) beliefs formed from private information with endogenous price formation, while retaining an equilibrium concept that is mathematically well-posed, computable, and connectable to data. This paper develops a mean field game (MFG) framework for a market populated by a continuum of investors who filter an unobserved fundamental from noisy private signals and trade against an endogenous price impacted by mean trading rate (order flow).

The model keeps a clean stock--flow separation. Inventory $x_{i,t}$ (shares) carries over between decision times, trading rate $u_{i,t}$ (shares per unit time) satisfies $x_{i,t+\Delta t}=x_{i,t}+u_{i,t}\Delta t$, and the price equation loads on the mean flow $\bar u_t$ rather than on mean inventory. Overconfidence enters only through \emph{filtering}: agents treat their private signal as more precise than it truly is, which changes posterior uncertainty and therefore trading intensity and disagreement, not through an ad hoc demand multiplier.

Early behavioral models introduced stylized agent types (noise traders, momentum traders) to demonstrate how bounded rationality can drive prices away from fundamentals \cite{DeLong1990,Hong1999,DanielHirshleiferSubrahmanyam1998,GervaisOdean2001}, but typically rely on a small number of types rather than a full distribution of beliefs. Mean field games have been applied to financial markets in several contexts: Gu\'eant, Lasry, and Lions \cite{GueantLasryLions2010} developed MFG frameworks for execution and portfolio problems; Carmona, Fouque, and Sun \cite{CarmonaFouqueSun2015} studied systemic risk; Lachapelle \emph{et al.} \cite{LachapelleLasryLehalleLions2016} applied MFG ideas to high-frequency trading and price formation; and Cardaliaguet and Lehalle \cite{CardaliaguetLehalle2018} analyzed trade crowding effects.

Our contribution relative to existing MFG finance literature is to provide a bounded-rational equilibrium concept that is both economically interpretable and mathematically checkable in a stock--flow microstructure setting. In the dynamic baseline, we study a discrete-time \emph{myopic mean-field equilibrium}: given a conjectured common-noise adapted mean trading rate $(\bar u_t)$, each agent applies a one-step \CARA best response under conditionally Gaussian price increments (using their filtered belief), and mean-field consistency closes the dynamics by requiring $\bar u_t$ to equal the conditional population mean of these best responses. Because the best response is linear in $(\hat v_{i,t}-p_t)$ and in $\bar u_t$, the closure reduces each time step to a scalar stock--flow fixed point. A per-step contraction condition guarantees uniqueness and well-posedness of the induced discrete-time dynamics and particle approximation (Theorem~\ref{thm:contraction}).

The simulations distinguish \emph{level amplification} (large mispricing/volatility levels driven by weak anchoring and high noise trading) from \emph{overconfidence effects} (incremental differences when increasing $k$). Across a targeted joint grid over anchoring, impact, and noise trading, we find a robust mechanism separation in the baseline anchored-impact regime: microstructure parameters govern mispricing levels, while overconfidence strongly increases disagreement and trading intensity with economically small effects on price-level mispricing (Subsection~\ref{subsec:joint_grid}). To make the model empirically usable, Section~\ref{sec:param_identification} provides an identification argument linking key parameters to observable price and volume moments, and a moment-match vignette illustrates the calibration bridge on a liquid ETF.

The paper's dynamic analysis has a deliberate two-model structure. (A) The object we compute and simulate is the discrete-time stock--flow myopic equilibrium described above. (B) As a theoretical benchmark, Subsection~\ref{subsec:ct_cara_lqg} derives an exact finite-horizon continuous-time \CARA/\LQG\ best response under partial information; this analytic extension clarifies how the full intertemporal solution adds a hedging correction to the myopic rule, and we compare the myopic and intertemporal policies in Subsection~\ref{subsec:extensions}.

The main contributions of this paper are:
\begin{itemize}
    \item We formulate a mean field game with overconfident Bayesian filtering and an explicit stochastic price equation with endogenous flow impact, under a stock--flow timing that separates inventory and trading rate (Subsubsection~\ref{subsubsec:timing_stock_flow}).
    \item We provide a clean static benchmark with a closed-form linear price \eqref{eq:static_price} and explicit comparative statics in the overconfidence parameter $k$.
    \item In the dynamic setting, we derive a checkable myopic \CARA inventory rule under conditionally Gaussian price increments (with holding cost) and obtain a per-step stock--flow mean-field closure; an explicit contraction condition guarantees uniqueness and discrete-time well-posedness (Theorem~\ref{thm:contraction}).
    \item We report particle-based numerical experiments (large-$N$ simulations) across bull/bear regimes and microstructure regimes, including mechanism diagnostics (belief dispersion and perceived uncertainty), validation checks, and extensions (Section~\ref{sec:numerical}).
    \item We provide a minimal identification/calibration bridge: Section~\ref{sec:param_identification} links key parameters to observable price and volume moments and illustrates the mapping with a moment-match vignette.
\end{itemize}

The remainder of this paper is organized as follows. Section~\ref{sec:model} presents the model, including the static benchmark and the dynamic price/belief dynamics. Section~\ref{sec:meanfield} discusses the mean-field object in the presence of common noise and the fixed-point notion of (myopic) equilibrium. Section~\ref{sec:results} states the main analytic results, including the myopic rule and the continuous-time \CARA/\LQG\ benchmark. Section~\ref{sec:proofs} collects proofs. Section~\ref{sec:numerical} presents numerical experiments, the empirical bridge, and the joint-grid analysis. Section~\ref{sec:conclusion} concludes.


\section*{Main contributions}
\begin{itemize}
    \item We formulate a mean field game with overconfident Bayesian filtering and an explicit stochastic price equation with endogenous flow impact, under a stock--flow timing that separates inventory and trading rate (Subsubsection~\ref{subsubsec:timing_stock_flow}).
    \item We provide a clean static benchmark with a closed-form linear price \eqref{eq:static_price} and explicit comparative statics in the overconfidence parameter $k$.
    \item In the dynamic setting, we derive a checkable myopic \CARA inventory rule under conditionally Gaussian price increments (with holding cost) and obtain a per-step stock--flow mean-field closure; an explicit contraction condition guarantees uniqueness and discrete-time well-posedness (Theorem~\ref{thm:contraction}).
    \item We report particle-based numerical experiments (large-$N$ simulations) across bull/bear regimes and microstructure regimes, including mechanism diagnostics (belief dispersion and perceived uncertainty), validation checks, and extensions (Section~\ref{sec:numerical}).
    \item We provide a minimal identification/calibration bridge: Section~\ref{sec:param_identification} links key parameters to observable price and volume moments and illustrates the mapping with a moment-match vignette.
\end{itemize}

% --------------------------------------------------------------------
% Technical content (preserving original order after the introduction).
% --------------------------------------------------------------------

\section{Problem Statement and Model}
\label{sec:model}
% Section 2: Problem Statement and Model

\subsection{Setup and Notation}

We consider a financial market with a continuum of investors indexed by $i \in [0,1]$ over a horizon $[0,T]$. Let $(\Omega,\mathcal{F},\{\mathcal{F}_t\}_{t\ge 0},\Prob)$ support independent Brownian motions $W_t$ (fundamental shock), $Z_t$ (aggregate price/noise-trading shock), $U_t$ (common signal component), and $\{B_{i,t}\}_{i\in[0,1]}$ (idiosyncratic signal noise). We write $\E[\cdot]$ for expectation and $\tau = 1/\sigma^2$ for precision.

Model parameters: $\gamma>0$ (absolute risk aversion), $\kappa>0$ (fundamental anchoring strength), $\lambda>0$ (impact strength), $k\ge 1$ (overconfidence), and volatilities $\sigma_v,\sigma_c,\sigma_\epsilon,\sigma_\eta,\phi>0$ (fundamental shock, common signal component, idiosyncratic signal noise, order-flow noise scale, and reduced-form price/noise-trading volatility). In the microstructure sketch below, $\sigma_\eta$ induces price noise via $\phi:=\lambda\sigma_\eta$. Overconfidence is modeled as a misperception of private-signal precision: an agent with parameter $k$ uses perceived idiosyncratic variance $\sigma_\epsilon^2/k$ in belief updating (equivalently, perceived precision $k/\sigma_\epsilon^2$).

\subsection{Static Benchmark}

In this static benchmark, $\nu\sim \mathcal{N}(0,\sigma_\nu^2)$ is an exogenous noise-demand shock (price level), distinct from the order-flow noise scale $\sigma_\eta$ used in the microstructure sketch (whose induced price-noise volatility is $\phi=\lambda\sigma_\eta$).

As a baseline, consider a static market with fundamental value $v \sim \mathcal{N}(\mu_{v,0},\sigma_{v,0}^2)$. Each investor receives a private signal
\begin{equation}
s_i = v + \epsilon_i, \quad \epsilon_i \sim \mathcal{N}(0, \sigma_{\epsilon,0}^2),
\end{equation}
and chooses demand $x_i$ for a single risky asset with payoff $v$ at settlement. Investors have \CARA utility
\begin{equation}
U_i = -\E\!\left[e^{-\gamma W_i}\mid s_i\right], \quad W_i = w_0 + x_i(v-p),
\end{equation}
and take the price impact rule (instantaneous clearing with noise demand)
\begin{equation}
p = \lambda \int_0^1 x_i\,di + \nu,\quad \nu \sim \mathcal{N}(0,\sigma_\nu^2).
\end{equation}
Under Gaussian beliefs, maximizing \CARA utility is equivalent to maximizing mean minus $(\gamma/2)$ times variance. Let $\tau_{v,0}:=1/\sigma_{v,0}^2$ and let $\tau'_{\epsilon,0}:=k/\sigma_{\epsilon,0}^2$ denote the perceived private-signal precision. Then the (perceived) posterior variance is $\sigma^2_{\text{post}} = 1/(\tau_{v,0}+\tau'_{\epsilon,0})$ and the posterior mean is $\hat v_i = (\tau_{v,0}\mu_{v,0} + \tau'_{\epsilon,0} s_i)/(\tau_{v,0}+\tau'_{\epsilon,0})$. The optimal demand is therefore
\begin{equation}
x_i^* = \frac{\hat v_i - p}{\gamma\,\sigma^2_{\text{post}}}.
\end{equation}
To make the aggregation step explicit, consider first a finite population of size $N$, with signals $s_i=v+\epsilon_i$ for $i=1,\dots,N$, where $\{\epsilon_i\}_{i=1}^N$ are i.i.d., independent of $v$, with $\E[\epsilon_i]=0$ and $\E[\epsilon_i^2]=\sigma_{\epsilon,0}^2<\infty$, and define the sample average $\bar s_N := \frac{1}{N}\sum_{i=1}^N s_i$.
Since $\hat v_i$ is affine in $s_i$, the cross-sectional mean belief satisfies
$
\frac{1}{N}\sum_{i=1}^N \hat v_i
=\frac{\tau_{v,0}\mu_{v,0}+\tau'_{\epsilon,0} \bar s_N}{\tau_{v,0}+\tau'_{\epsilon,0}}.
$
Under the finite-$N$ clearing rule $p^N=\lambda\big(\frac1N\sum_{i=1}^N x_i\big)+\nu$, the same algebra yields
\[
p^N
= \frac{\lambda \tau'_{\epsilon,0}}{\gamma + \lambda(\tau_{v,0}+\tau'_{\epsilon,0})}\,\bar s_N
 + \frac{\gamma}{\gamma + \lambda(\tau_{v,0}+\tau'_{\epsilon,0})}\,\nu
 + \frac{\lambda\tau_{v,0}}{\gamma + \lambda(\tau_{v,0}+\tau'_{\epsilon,0})}\,\mu_{v,0}.
\]
Since $\bar s_N=v+\frac1N\sum_{i=1}^N \epsilon_i$, the strong law of large numbers implies $\bar s_N\to v$ almost surely (conditional on $v$, hence also unconditionally) as $N\to\infty$, yielding the mean-field limit price\footnote{A literal continuum model with $\int_0^1 s_i\,di=v$ can be formalized under an ``exact law of large numbers'' framework for a continuum of independent shocks (e.g., via a Fubini extension/Loeb space). We avoid this measure-theoretic detour and adopt the large-$N$ interpretation, consistent with the finite-$N$ particle/limit perspective used later.}
\begin{equation}
\begin{aligned}
p
&= \frac{\lambda \tau'_{\epsilon,0}}{\gamma + \lambda(\tau_{v,0}+\tau'_{\epsilon,0})}\,v
 + \frac{\gamma}{\gamma + \lambda(\tau_{v,0}+\tau'_{\epsilon,0})}\,\nu\\
&\quad + \frac{\lambda\tau_{v,0}}{\gamma + \lambda(\tau_{v,0}+\tau'_{\epsilon,0})}\,\mu_{v,0}.
\end{aligned}
\label{eq:static_price}
\end{equation}
Equation \eqref{eq:static_price} makes the channel explicit: higher $k$ increases perceived precision $\tau'_{\epsilon,0}$ and therefore strengthens the price response to $v$ while also increasing demand magnitude for a given mispricing.

\paragraph{Notation note.}
In the static benchmark (Sec.~2.2), $(\mu_{v,0},\sigma_{v,0})$ parameterize the one-shot prior $v\sim\mathcal N(\mu_{v,0},\sigma_{v,0}^2)$ and $\sigma_{\epsilon,0}$ is the one-shot private-signal noise standard deviation in $s_i=v+\epsilon_i$. We retain the ``,0'' subscripts throughout Section~2.2 and its proofs; unsubscripted $(\mu_v,\sigma_v,\sigma_\epsilon,\ldots)$ are reserved for the dynamic model. In the dynamic model (Sec.~2.3), $(\mu_v,\sigma_v)$ denote the drift and instantaneous diffusion volatility of $v_t$ in $dv_t=\mu_v\,dt+\sigma_v\,dW_t$, and $\sigma_\epsilon$ is the diffusion loading on idiosyncratic signal noise in $d\xi_{i,t}=v_t\,dt+\sigma_c\,dU_t+\sigma_\epsilon\,dB_{i,t}$. The benchmarks are comparable but the parameters have different time-aggregation interpretations.

\subsection{Dynamic Mean Field Game Formulation}

\subsubsection{State Dynamics and Mean-Field Coupling}

The fundamental value follows a diffusion
\begin{equation}
dv_t = \mu_v\,dt + \sigma_v\,dW_t,
\end{equation}
and each investor observes a private signal process
 \begin{equation}
 d\xi_{i,t} = v_t dt + \sigma_c\,dU_t + \sigma_\epsilon\,dB_{i,t},
 \end{equation}
 where $\{B_{i,t}\}$ are independent across $i$.
 In the discrete-time implementation on a grid with step size $\Delta t$, we work with the normalized signal increment (we reserve $y_{i,t}$ below for mispricing $\hat v_{i,t}-p_t$):
 \[
 z_{i,t}:=\frac{\xi_{i,t+\Delta t}-\xi_{i,t}}{\Delta t}
 = v_t+\frac{\sigma_c}{\sqrt{\Delta t}}\,\varepsilon^c_t+\frac{\sigma_\epsilon}{\sqrt{\Delta t}}\,\varepsilon^i_{i,t},
 \qquad \varepsilon^c_t,\varepsilon^i_{i,t}\sim\mathcal N(0,1),
 \]
 so $\Var(z_{i,t}\mid v_t)=(\sigma_c^2+\sigma_\epsilon^2)/\Delta t$ and the perceived measurement variance is
 $R'_i(\Delta t)=(\sigma_c^2+\sigma_\epsilon^2/k_i)/\Delta t$ (Appendix~A). (For $\Delta t=1$, this reduces to
 $z_{i,t}=v_t+\sigma_c\varepsilon^c_t+\sigma_\epsilon\varepsilon^i_{i,t}$.)
The market price evolves endogenously through an anchored impact rule:
\begin{equation}
dp_t = \kappa (v_t - p_t)\,dt + \lambda \bar u_t\,dt + \phi\,dZ_t,\qquad \bar u_t := \int_0^1 u_{i,t}\,di.
\label{eq:anchored_impact}
\end{equation}
Equation \eqref{eq:anchored_impact} is a reduced-form microstructure: prices mean-revert to fundamentals at rate $\kappa$ (arbitrage/fundamental anchoring), respond to aggregate net trading flow with strength $\lambda$, and are perturbed by aggregate noise $Z_t$ with volatility $\phi$.
The anchoring term $\kappa(v_t-p_t)$ should be interpreted as reduced-form arbitrage pressure toward the (latent) fundamental $v_t$; it is not assumed that a market maker directly observes $v_t$.
The mean-field coupling enters through $\bar u_t$.
In the discrete-time implementation, $u_{i,t}$ is the trading rate induced by inventory adjustments over a decision interval $\Delta t$ (Subsubsection~\ref{subsubsec:timing_stock_flow}).

\paragraph{Microstructure sketch for the anchored impact rule.}
Equation~\eqref{eq:anchored_impact} can be viewed as a reduced-form limit of competitive
risk-neutral market making with inventory costs, plus an (exogenous) arbitrage/anchoring flow.

\emph{Linear impact from inventory-averse competitive market makers.}
Let $Q_t$ denote aggregate net order flow hitting dealers (positive when the public buys),
and let the representative dealer's inventory be $q_t$, with
\[
dq_t = -\,dQ_t, \qquad \text{(dealers take the opposite side of net flow).}
\]
Assume dealers are risk-neutral and competitive (zero expected profit), but face a quadratic
inventory holding cost $\frac{\lambda}{2}q_t^2$.
A standard myopic quoting / LQ argument implies the midprice is shifted linearly by inventory,
\[
p_t \approx \tilde m_t - \lambda q_t,
\]
where $\tilde m_t$ is the dealer's ``efficient'' benchmark (e.g., a fast-moving reference level).
Differentiating and substituting $dq_t=-dQ_t$ gives
\[
dp_t \approx d\tilde m_t + \lambda\, dQ_t,
\]
so that (in drift form) price changes are \emph{linear} in contemporaneous net flow. Writing
$dQ_t = \bar u_t\,dt + \sigma_\eta\,dZ_t$ (where $\bar u_t$ is the net order-flow intensity hitting dealers) yields the impact and noise terms
\[
dp_t-d\tilde m_t \approx \lambda \bar u_t\,dt + (\lambda\sigma_\eta)\, dZ_t,
\qquad \phi := \lambda\sigma_\eta.
\]
Here $\tilde m_t$ is the dealer's efficient benchmark, $\bar u_t$ is the net order-flow intensity hitting dealers, and $Z_t$ is a Brownian motion; $\sigma_\eta$ is the corresponding order-flow noise scale.
In the baseline below we explicitly distinguish inventory $x_{i,t}$ from trading rate $u_{i,t}$ (Subsubsection~\ref{subsubsec:timing_stock_flow}); impact in \eqref{eq:anchored_impact} depends on the mean trading rate $\bar u_t$ (equivalently $\lambda\,\Delta\bar x_t$ over a decision interval).

\emph{Anchoring as reduced-form arbitrage pressure.}
To capture fast arbitrage / fundamental-anchoring, posit an additional ``arbitrage'' flow that
leans against mispricing,
\[
u^{\mathrm{arb}}_t = a\,(v_t - p_t),
\]
so total net flow is $\bar u_t + u^{\mathrm{arb}}_t$ at the same impact slope $\lambda$.
Then
\[
dp_t \approx \lambda\big(\bar u_t + a(v_t-p_t)\big)\,dt + (\lambda\sigma_\eta)\,dZ_t
= \kappa(v_t-p_t)\,dt + \lambda \bar u_t\,dt + \phi\,dZ_t,
\quad \kappa := \lambda a,
\]
which matches the anchored impact specification.

\subsubsection{Equilibrium Price Formation}

We focus on feedback strategies where demand depends on perceived mispricing. Let $\hat v_{i,t}$ denote investor $i$'s Kalman--Bucy estimate of $v_t$ under the \emph{perceived} measurement-noise variance $R'(k)$ (Appendix~A). When $k=1$ this coincides with the true conditional mean $\E[v_t\mid \mathcal F^i_t]$; when $k\ne 1$ it should be interpreted as a \emph{subjective} filter state computed under the investor's misperceived signal precision. Define the mean belief $m_t := \int_0^1 \hat v_{i,t}\,di$. A broad class of (best-response and learning-based) specifications take the form
\begin{equation}
x_{i,t} = \Theta_t(\Sigma_{i,t})\bigl(\hat v_{i,t}-p_t\bigr),
\label{eq:lin_policy}
\end{equation}
where $\Sigma_{i,t}$ is the agent's perceived posterior variance.
\emph{Definition (gain).} For each $t$, $\Theta_t:(0,\infty)\to[0,\infty)$ maps perceived posterior variance to the (shares/\$) sensitivity of demand to mispricing and is weakly decreasing in $\Sigma$ (equivalently increasing in precision $1/\Sigma$). In the myopic \CARA best response, $\Theta_t$ is derived explicitly in Eq.~\eqref{eq:theta_def} and is nonnegative under $\gamma>0$, $\chi\ge0$, and $\sigma^2_{p,i,t}(\Sigma)>0$.

\subsubsection{Timing and stock--flow separation (inventory vs.\ trading rate)}
\label{subsubsec:timing_stock_flow}

\paragraph{Timing and objects (discrete decision interval $\Delta t$).}
At each decision date $t$ agents enter with a \emph{pre-trade inventory} $x_{i,t^-}$ (shares) and choose a \emph{post-trade inventory} $x_{i,t}$ (shares). The executed order and trading rate are
\[
\Delta x_{i,t} := x_{i,t}-x_{i,t^-},\qquad
u_{i,t} := \frac{\Delta x_{i,t}}{\Delta t}\quad (\text{shares/time}),\qquad
\bar u_t := \int_0^1 u_{i,t}\,di.
\]
Price changes over $[t,t+\Delta t]$ via \emph{flow-based} impact,
\[
\Delta p_{t+\Delta t}
= \kappa (v_t-p_t)\,\Delta t + \lambda\,\bar u_t\,\Delta t + \phi\sqrt{\Delta t}\,\varepsilon_{t+\Delta t},
\qquad \varepsilon_{t+\Delta t}\sim\mathcal N(0,1).
\]
Equivalently, since $\bar u_t\Delta t = \Delta\bar x_t$ with $\bar x_t:=\int_0^1 x_{i,t}\,di$, the impact term can be written $\lambda\,\Delta\bar x_t$.
This timing makes explicit that $x$ is a stock (shares) while $u$ is a flow (shares/time), resolving the unit ambiguity when a single symbol is used for both.

\subsubsection{Belief Updating via Kalman Filtering}

Each investor maintains a belief state $(\hat v_{i,t},\Sigma_{i,t})$ where $\hat v_{i,t}$ is the posterior mean of $v_t$ and $\Sigma_{i,t}$ is the perceived posterior variance. In the linear-Gaussian setting above, $(\hat v_{i,t},\Sigma_{i,t})$ satisfy Kalman--Bucy equations driven by the private signal $\xi_{i,t}$. Overconfidence enters by replacing $\sigma_\epsilon^2$ with $\sigma_\epsilon^2/k$ in the perceived filter. We summarize the resulting filter in Appendix~A.

\subsubsection{Objective Functional}

Investor $i$ chooses a post-trade inventory process $(x_{i,t})_{t}$ (shares) on the discrete decision grid.
Over one step, wealth evolves as
\[
W_{i,t+\Delta t} = W_{i,t} + x_{i,t}\,\Delta p_{t+\Delta t} - \frac{\chi}{2}x_{i,t}^2\,\Delta t,
\qquad \chi>0,
\]
where the quadratic term penalizes inventory holding/risk over the interval.
Rather than solving the full intertemporal exponential-utility control problem, we adopt a bounded-rational \emph{myopic} approximation: at each decision date, the investor solves the one-step certainty-equivalent problem
\[
\max_{x_{i,t}}\;
x_{i,t}\,\mu_{i,t}
- \frac12\big(\gamma\,\sigma^2_{p,i,t}+\chi\big)x_{i,t}^2\,\Delta t,
\]
where
$\mu_{i,t}:=\E_i[\Delta p_{t+\Delta t}\mid\mathcal F^i_t]=\big(\kappa(\hat v_{i,t}-p_t)+\lambda\bar u_t\big)\Delta t$
and $\Var_i(\Delta p_{t+\Delta t}\mid\mathcal F^i_t)=\sigma^2_{p,i,t}\Delta t$.
We treat $\sigma^2_{p,i,t}$ as computed from primitives; in the discrete implementation it includes belief uncertainty through the effective variance term in Section~IV.

\subsubsection{Mean-Field Consistency (Myopic)}

Let $\mu_t$ denote the (possibly random) law of $(\hat v_{i,t},\Sigma_{i,t})$ across agents. Together with \eqref{eq:anchored_impact}, we study a \emph{myopic mean-field consistent} fixed point: (i) given a conjectured mean-field flow $\{\mu_t\}$ (equivalently, moments such as $m_t$), each agent applies the myopic \CARA rule in Section~IV; and (ii) the resulting population law induced by the controlled belief dynamics is consistent with $\{\mu_t\}$. Section~III formalizes the mean-field object and discusses the common-noise aspect induced by $W_t$ and $Z_t$.


\section{Mean-Field Approximation and Limiting Equations}
\label{sec:meanfield}
% Section 3: Mean-Field Approximation and Limiting Equations

\subsection{Mean-Field Object Under Common Noise}

Because the model includes common shocks $(W_t,U_t,Z_t)$ through the fundamental, correlated signals, and price noise, the appropriate mean-field object is the \emph{conditional} population law given the common filtration $\mathcal{F}^0_t := \sigma(W_s,U_s,Z_s: s\le t)$. Concretely, let $X_{i,t}:=(\hat v_{i,t},\Sigma_{i,t},k_i)$ denote investor $i$'s internal state (belief, perceived uncertainty, and overconfidence type), and define the mean-field flow
\begin{equation}
\mu_t := \Law(X_{i,t}\mid \mathcal{F}^0_t).
\end{equation}
The conditional mean belief $m_t := \int \hat v\,\mu_t(d\hat v\,d\Sigma\,dk)$ enters the price dynamics through the mean trading rate $\bar u_t$ appearing in \eqref{eq:anchored_impact}; inventories and trading rates are linked by the stock--flow timing in Subsubsection~\ref{subsubsec:timing_stock_flow}.

\subsection{Mean-Field Equilibrium (Common-Noise Formulation)}

Fix a candidate flow $\{\mu_t\}_{t\in[0,T]}$ and the resulting price process $\{p_t\}$ from \eqref{eq:anchored_impact}. In the bounded-rational formulation of this paper, the representative agent does not solve a full intertemporal control problem; instead, given $\{\mu_t\}$ and $\{p_t\}$ the agent applies the myopic rule from Section~IV, yielding a feedback $x_{i,t} = x^{\mathrm{myopic}}(t,X_{i,t},p_t,\mu_t)$. The induced controlled state process generates a new conditional flow $\{\tilde\mu_t\}$. A myopic mean-field equilibrium is a fixed point $\mu=\tilde\mu$.

In common-noise settings, $\mu_t$ is random and its evolution is governed by a stochastic Kolmogorov forward equation (stochastic Fokker--Planck / SPDE) for the conditional law; see, e.g., \cite{CarmonaDelarueMFGVol1,Lacker2016}. In Section~VI we compute equilibrium statistics using a particle approximation: simulate a large $N$-agent system driven by shared $(W,Z)$ and independent $\{B_i\}$ and estimate $\mu_t$ via the empirical measure.

\paragraph{Scope of theoretical results.}
To avoid confusion, we explicitly state what this paper establishes and what it does not claim.
\textbf{Baseline object (what we compute and simulate):} a discrete-time, stock--flow myopic equilibrium with inventory carryover and flow-based impact. The per-step mean-field fixed point is scalar and has a unique solution when $0\le B_t<\Delta t$ (Theorem~\ref{thm:contraction}); under this condition the full discrete-time particle system is well-posed (Theorem~\ref{thm:well_posed}).
\textbf{Benchmarks (for comparison only):} we also record a continuous-time position-impact benchmark (which collapses stock and flow) where the mean-field closure becomes algebraic and admits an explicit closed form and condition (Theorem~\ref{thm:cn_fp}); this benchmark is not used in the baseline simulations.
\textbf{What we do not claim:} This paper does not establish existence or uniqueness of a full optimal-control MFG equilibrium under intertemporal exponential utility maximization with endogenous filtering on prices. The myopic rule \eqref{eq:myopic_policy} is a bounded-rational approximation; Subsection~\ref{subsec:ct_cara_lqg} and Subsection~\ref{subsec:ct_flow_control_extension} provide constructive continuous-time comparisons, but a complete verification of existence/uniqueness conditions for those intertemporal MFG couplings would strengthen the theoretical foundation.

\subsection{Discrete-Time Myopic Mean-Field Equilibrium (What We Compute)}

The numerical experiments in Section~VI implement a discrete-time, bounded-rational equilibrium concept consistent with the myopic rule in Section~IV. Fix a time grid $t=0,1,\dots,T$ (step size $\Delta t=1$ in the implementation) and consider the $N$-agent system induced by:
(i) fundamental and price updates (Euler discretization of \eqref{eq:anchored_impact}),
(ii) private-signal observations (normalized increments $z_{i,t}=(\xi_{i,t+\Delta t}-\xi_{i,t})/\Delta t$) and Kalman belief updates (Appendix~A), and
(iii) the myopic inventory rule \eqref{eq:myopic_policy} with the per-step stock--flow fixed point (Proposition~\ref{prop:step_fixed_point}).

\textbf{Definition (myopic mean-field equilibrium, discrete time).}
Given a pre-trade mean inventory $\bar x_{t^-}$ and public price $p_t$, agent $i$ computes a belief state $(\hat v_{i,t},\Sigma_{i,t})$ and chooses a post-trade inventory $x_{i,t}$ according to the myopic rule, which depends on the mean trading rate $\bar u_t$. A discrete-time myopic mean-field equilibrium at time $t$ is a pair $(\bar x_t,\bar u_t)$ such that
\[
\bar x_t = \frac{1}{N}\sum_{i=1}^N x_{i,t},
\qquad
\bar u_t=\frac{\bar x_t-\bar x_{t^-}}{\Delta t},
\]
with individual choices consistent with $\bar u_t$. In the implementation the rule is linear in $\bar u_t$ and the scalar fixed point is computed in closed form (Proposition~\ref{prop:step_fixed_point}).

\subsection{Discrete-Time Well-Posedness (Existence and Uniqueness of Paths)}

We formalize the particle implementation as a coupled discrete-time dynamical system on a fixed time grid.

\textbf{Assumptions.}
Fix $T\in\mathbb N$ and $\Delta t>0$ and assume:
(A1) parameters satisfy $\kappa>0$, $\phi>0$, $\gamma>0$, $\alpha_0>0$ and initial conditions $(v_0,p_0,\{\hat v_{i,0},\Sigma_{i,0}\}_{i=1}^N)$ are finite with $\Sigma_{i,0}\ge 0$;
(A2) beliefs evolve via the discrete-time Kalman recursion in Appendix~A (so $(\hat v_{i,t},\Sigma_{i,t})$ is uniquely defined given $(\hat v_{i,t-1},\Sigma_{i,t-1})$ and the signal);
(A3) for each $t$, the per-step stock--flow map admits a unique fixed point $\bar x_t$ (and hence $\bar u_t$), e.g., the contraction condition $0\le B_t<\Delta t$ in Theorem~\ref{thm:contraction} holds;
(A4) price and fundamental updates are given by the Euler discretizations of their SDEs driven by specified shock sequences.

\begin{theorem}[Well-posedness of the discrete-time particle system]
\label{thm:well_posed}
Under (A1)--(A4), for any realization of the common and idiosyncratic shock sequences, the discrete-time system defined by: (i) belief update (Appendix~A), (ii) per-step myopic mean-field fixed point (Proposition~\ref{prop:step_fixed_point}), and (iii) Euler updates for $(v_t,p_t)$, admits a unique adapted solution path
\[
(v_t,p_t,\{\hat v_{i,t},\Sigma_{i,t}\}_{i=1}^N,\bar x_t,\bar u_t)_{t=0}^T.
\]
\end{theorem}

\subsection{Position-impact benchmark: common-noise fixed-point verification (continuous time)}\label{sec:cn_fp_verification}

This section is a diagnostic benchmark using the legacy position-impact closure; the baseline simulations use the discrete-time stock--flow equilibrium with uniqueness condition $0\le B_t<\Delta t$ (Proposition~\ref{prop:step_fixed_point}, Theorem~\ref{thm:contraction}).

This subsection records a self-contained fixed-point verification for a continuous-time \emph{position-impact} benchmark (which collapses stock and flow so that $\bar u_t\equiv \bar x_t$).
It is included for comparison and is not used in the baseline simulations, which use the discrete-time stock--flow closure above.
The key point is that, in our information simplification (filtering on $\xi^i$ only),
the belief state dynamics are \emph{conditionally i.i.d.} given the common filtration, so the mean-field
object is the conditional law and the equilibrium closure reduces to an explicit algebraic fixed point.

\paragraph{Common and idiosyncratic noises.}
Recall the common filtration
\[
\mathcal F^0_t := \sigma(W_s,U_s,Z_s:\; s\le t),
\]
and idiosyncratic signal noises $\{B^i\}_{i\ge 1}$ independent across $i$ and independent of $(W,U,Z)$.

\paragraph{Finite-$N$ (continuous-time) price-taking system.}
Consider the $N$-agent system (continuous time) with a \emph{price-taking} interaction through the
average demand $\bar x^N_t := \frac1N\sum_{j=1}^N x^{j,N}_t$:
\begin{align}
dv_t &= \mu_v\,dt + \sigma_v\,dW_t, \label{eq:cn_v}\\
d\xi^i_t &= v_t\,dt + \sigma_c\,dU_t + \sigma_\epsilon\,dB^i_t,\qquad i=1,\dots,N, \label{eq:cn_signal}\\
dp^N_t &= \kappa(v_t-p^N_t)\,dt + \lambda \bar x^N_t\,dt + \phi\,dZ_t. \label{eq:cn_priceN}
\end{align}
Agents compute the subjective Kalman--Bucy filter (Appendix~A) from \eqref{eq:cn_signal}:
\begin{equation}\label{eq:cn_filter}
\begin{split}
d\hat v^{i}_t &= \mu_v\,dt + K_t^{(i)}\bigl(d\xi^i_t - \hat v^{i}_t\,dt\bigr),\\
K_t^{(i)} &:= \frac{\Sigma_t^{(i)}}{R^{\prime}_i},\quad R^{\prime}_i:=\sigma_c^2+\sigma_\epsilon^2/k_i.
\end{split}
\end{equation}
with Riccati $\dot\Sigma_t^{(i)}=\sigma_v^2-(\Sigma_t^{(i)})^2/R_i'$ and $\Sigma_0^{(i)}\ge 0$.
(For a homogeneous population, $k_i\equiv k$ and hence $K_t^{(i)}\equiv K_t$ and $\Sigma_t^{(i)}\equiv \Sigma_t$
are deterministic.)

Given $(p^N_t,\bar x^N_t)$, the (benchmark) myopic CARA rule under Gaussian price increments (with $\chi=0$ and $\Var(dp_t\mid\mathcal F_t^i)/dt=\phi^2$) prescribes the pointwise best response
\begin{equation}\label{eq:cn_myopic_BR}
x^{i,N}_t
= \frac{\kappa(\hat v^i_t-p^N_t)+\lambda \bar x^N_t}{\gamma\phi^2},
\qquad i=1,\dots,N.
\end{equation}
Averaging \eqref{eq:cn_myopic_BR} yields the \emph{per-time} fixed point
\begin{equation}\label{eq:cn_fpN}
(\gamma\phi^2-\lambda)\,\bar x^N_t = \kappa(\hat m^N_t-p^N_t),
\qquad
\hat m^N_t:=\frac1N\sum_{i=1}^N \hat v^i_t,
\end{equation}
so whenever $\gamma\phi^2>\lambda$,
\begin{equation}\label{eq:cn_barxN_closed}
\bar x^N_t = \frac{\kappa}{\gamma\phi^2-\lambda}\,(\hat m^N_t-p^N_t).
\end{equation}

\paragraph{Mean-field (conditional-law) system.}
Define the representative-agent belief state $\hat v_t$ solving \eqref{eq:cn_filter} with an independent copy
$B$ (and the same common noises $(W,U,Z)$), and define the conditional flow
\[
\begin{aligned}
\mu_t &:= \mathcal L\bigl(\hat v_t,\Sigma_t,k \,\big|\, \mathcal F^0_t\bigr),\\
m_t &:= \int \hat v\,\mu_t(d\hat v\,d\Sigma\,dk)=\mathbb E[\hat v_t\mid \mathcal F^0_t].
\end{aligned}
\]
The mean-field price satisfies
\begin{equation}\label{eq:cn_price_limit}
dp_t = \kappa(v_t-p_t)\,dt + \lambda \bar x_t\,dt + \phi\,dZ_t,
\end{equation}
where $\bar x_t$ is required to satisfy the mean-field consistency condition
\begin{equation}\label{eq:cn_mf_consistency}
\bar x_t
= \int \frac{\kappa(\hat v-p_t)+\lambda \bar x_t}{\gamma\phi^2}\ \mu_t(d\hat v\,d\Sigma\,dk).
\end{equation}
Since only the mean of $\hat v$ enters \eqref{eq:cn_mf_consistency}, the fixed point reduces to
\begin{equation}\label{eq:cn_barx_closed}
\bar x_t = \frac{\kappa}{\gamma\phi^2-\lambda}\,(m_t-p_t), \qquad \text{assuming }\gamma\phi^2>\lambda.
\end{equation}

\begin{theorem}[Common-noise fixed point]\label{thm:cn_fp}
Assume $\gamma\phi^2>\lambda$ and $\sup_{i\le N}\mathbb E[|\hat v^i_0|^2]<\infty$ with
$\sup_{i\le N}\Sigma^{(i)}_0<\infty$.
For simplicity, assume initial belief states are i.i.d.\ across $i$ conditional on $\mathcal F^0_0$.
Then:

\smallskip
\noindent\textup{(i) (Well-posedness.)} The $N$-agent system \eqref{eq:cn_v}--\eqref{eq:cn_priceN} with
controls given by the myopic fixed point \eqref{eq:cn_myopic_BR} admits a unique strong solution, and the limiting
mean-field (conditional-law) system consisting of \eqref{eq:cn_v}, a representative signal equation of the form \eqref{eq:cn_signal} (with an independent idiosyncratic Brownian motion), the representative filter \eqref{eq:cn_filter}, and the price equation \eqref{eq:cn_price_limit} with $\bar x$ given by \eqref{eq:cn_barx_closed} admits a unique strong solution.

\smallskip
\noindent\textup{(ii) (Fixed-point verification.)} The process $\bar x$ defined by \eqref{eq:cn_barx_closed} is the unique solution
to the mean-field consistency equation \eqref{eq:cn_mf_consistency} (i.e., a unique myopic common-noise MFG fixed point).

\smallskip
\noindent\textup{(iii) (Conditional propagation of chaos / LLN.)} Let the empirical measure of belief states be
$\mu^N_t := \frac1N\sum_{i=1}^N \delta_{(\hat v^i_t,\Sigma^{(i)}_t,k_i)}$. Then for each fixed $t$,
\[
\mu^N_t \xrightarrow[N\to\infty]{} \mu_t
\quad\text{in probability, and}\quad
\hat m^N_t \to m_t \ \text{in }L^2.
\]
Consequently, $\bar x^N_t\to \bar x_t$ in $L^2$ and $p^N_t\to p_t$ in $L^2$ uniformly on $[0,T]$.

\end{theorem}

\begin{proof}
\emph{Step 1 (Well-posedness).}
For each $i$, the filter SDE \eqref{eq:cn_filter} is linear with adapted coefficients and admits a unique strong solution
with finite second moments; the Riccati equation is a scalar ODE with global solution (Appendix~A).
Given $\bar x^N$ and $p^N$, the control \eqref{eq:cn_myopic_BR} is affine in $(\hat v^i,p^N,\bar x^N)$.
The averaging identity \eqref{eq:cn_fpN} yields the explicit closure \eqref{eq:cn_barxN_closed} when
$\gamma\phi^2>\lambda$. Substituting \eqref{eq:cn_barxN_closed} into \eqref{eq:cn_priceN} gives a linear SDE for $p^N$
driven by $(W,U,Z)$ and the averaged belief $\hat m^N$, hence a unique strong solution. Likewise, substituting
\eqref{eq:cn_barx_closed} into \eqref{eq:cn_price_limit} gives a linear SDE for $p$ driven by $(W,U,Z)$ and $m$.

\emph{Step 2 (Fixed point).}
Equation \eqref{eq:cn_mf_consistency} is affine in $\bar x_t$:
\[
\bar x_t = \frac{\kappa(m_t-p_t)+\lambda \bar x_t}{\gamma\phi^2}.
\]
Rearranging gives $(\gamma\phi^2-\lambda)\bar x_t=\kappa(m_t-p_t)$ and thus \eqref{eq:cn_barx_closed}.
Uniqueness follows since $\gamma\phi^2-\lambda>0$.

\emph{Step 3 (Conditional LLN / propagation of chaos).}
Conditional on $\mathcal F^0_t$, the idiosyncratic noises $\{B^i\}$ are independent and the filter equations
\eqref{eq:cn_filter} are driven by independent Brownian motions with the \emph{same} common coefficients and common path
$(W,U,Z)$. Hence the collection $\{(\hat v^i_t,\Sigma^{(i)}_t,k_i)\}_{i=1}^N$ is i.i.d.\ conditional on $\mathcal F^0_t$
(with law $\mu_t$), and the conditional law of large numbers implies $\mu^N_t\Rightarrow \mu_t$ and
\[
 \mathbb E\bigl[|\hat m^N_t-m_t|^2\mid \mathcal F^0_t\bigr]
 = \frac{1}{N}\,\mathrm{Var}(\hat v_t\mid \mathcal F^0_t),
 \]
 so $\hat m^N_t\to m_t$ in $L^2$. Moreover, by Step 1 we have $\sup_{t\le T}\mathbb E[\hat v_t^2]<\infty$, hence
 $\sup_{t\le T}\mathbb E[\mathrm{Var}(\hat v_t\mid \mathcal F^0_t)]\le \sup_{t\le T}\mathbb E[\hat v_t^2]<\infty$ and therefore
 \[
 \sup_{t\le T}\mathbb E|\hat m^N_t-m_t|^2 \le \frac{C_T}{N},
 \qquad
 \int_0^T \mathbb E|\hat m^N_s-m_s|^2\,ds \le \frac{TC_T}{N}\xrightarrow[N\to\infty]{}0.
 \]

\emph{Step 4 (Price convergence).}
Let $C:=\kappa/(\gamma\phi^2-\lambda)$ so that \eqref{eq:cn_barxN_closed}--\eqref{eq:cn_barx_closed} read
$\bar x^N_t=C(\hat m^N_t-p^N_t)$ and $\bar x_t=C(m_t-p_t)$. Substituting these into
\eqref{eq:cn_priceN} and \eqref{eq:cn_price_limit} gives the closed price equations
\[
dp^N_t = \kappa(v_t-p^N_t)\,dt + \lambda C(\hat m^N_t-p^N_t)\,dt + \phi\,dZ_t,
\qquad
dp_t = \kappa(v_t-p_t)\,dt + \lambda C(m_t-p_t)\,dt + \phi\,dZ_t.
\]
Subtracting yields a linear ODE for $\Delta p_t:=p^N_t-p_t$ driven by $\Delta m_t:=\hat m^N_t-m_t$:
\[
d\Delta p_t = -(\kappa+\lambda C)\Delta p_t\,dt + \lambda C\,\Delta m_t\,dt.
\]
By variation of constants, for $t\in[0,T]$,
\[
\Delta p_t = e^{-(\kappa+\lambda C)t}\Delta p_0 + \lambda C\int_0^t e^{-(\kappa+\lambda C)(t-s)}\Delta m_s\,ds.
\]
Assuming matching initial conditions ($\Delta p_0=0$), Jensen and Cauchy--Schwarz imply
\[
\sup_{t\le T}\mathbb E|p^N_t-p_t|^2
 \ \le\ (\lambda C)^2 T \int_0^T \mathbb E|\hat m^N_s-m_s|^2\,ds
 \ \xrightarrow[N\to\infty]{}\ 0,
 \]
and the $L^2$ uniform convergence follows.

\emph{Step 5 (Mean-field action convergence).}
Using again $\bar x^N_t=C(\hat m^N_t-p^N_t)$ and $\bar x_t=C(m_t-p_t)$,
\[
\bar x^N_t-\bar x_t
=C\bigl[(\hat m^N_t-m_t)-(p^N_t-p_t)\bigr].
\]
Thus, by Step 3 and Step 4,
\[
\sup_{t\le T}\mathbb E|\bar x^N_t-\bar x_t|^2
\le 2C^2\left(\sup_{t\le T}\mathbb E|\hat m^N_t-m_t|^2+\sup_{t\le T}\mathbb E|p^N_t-p_t|^2\right)
\xrightarrow[N\to\infty]{}0.
\]
\end{proof}

\begin{remark}[Stochastic Fokker--Planck equation for $\mu_t$]\label{rem:cn_spde}
For any test function $\varphi\in C_b^2(\mathbb R\times \mathbb R_+\times \mathbb R_+)$,
the conditional law $\mu_t=\mathcal L(\hat v_t,\Sigma_t,k\mid \mathcal F^0_t)$ satisfies the common-noise SPDE
(in weak form)
\[
\begin{aligned}
d\langle \mu_t,\varphi\rangle
&= \Big\langle \mu_t,\bigl(\mu_v + K_t(k)(v_t-\hat v)\bigr)\,\partial_{\hat v}\varphi\\
&\qquad +\tfrac{1}{2} K_t(k)^2\sigma_\epsilon^2\,\partial^2_{\hat v\hat v}\varphi\Big\rangle dt\\
&\quad +\Big\langle \mu_t, K_t(k)\sigma_c\,\partial_{\hat v}\varphi\Big\rangle dU_t.
\end{aligned}
\]
where $\langle\mu,\varphi\rangle:=\int \varphi\,d\mu$ and
$K_t(k)=\Sigma_t(k)/(\sigma_c^2+\sigma_\epsilon^2/k)$.
In the homogeneous case, $\mu_t$ is Gaussian in $\hat v$ conditional on $\mathcal F^0_t$ (with random mean $m_t$).
\end{remark}

% Mean-Field Limit (black-box framework)
% To be included where the main text references ``Mean-Field Limit''.
% Uses \cite{CarmonaDelarueMFGVol1}, \cite{CarmonaDelarueMFGVol2}, \cite{Lacker2018},
% \cite{FournierGuillin2015}, \cite{HuangMalhameCaines2006} --- keys to be added in refs.bib.

\subsection{Mean-Field Limit}\label{sec:mf_limit}

We state a standard common-noise McKean--Vlasov/MFG limit as a black-box framework; our contribution is the constructive solution in the LQG specialization, given in Section~\ref{sec:lqg_constructive}.

\paragraph{Setting.}
Consider a common-noise McKean--Vlasov system on a finite horizon $[0,T]$. Let $\mathcal{F}^0_t := \sigma(W^0_s: s\le t)$ be the filtration generated by a common Brownian motion $W^0$, and for each $N\ge 1$ let $\{B^i\}_{i=1}^N$ be mutually independent and independent of $W^0$. The $N$-agent state processes $\{X^{i,N}_t\}_{i=1}^N$ take values in $\R^d$ and satisfy a coupled system whose coefficients may depend on the empirical measure $\bar\mu^N_t := \frac{1}{N}\sum_{j=1}^N \delta_{X^{j,N}_t}$ and on the common noise. The mean-field limit is described by a McKean--Vlasov SDE in which the conditional law $\mu_t = \Law(X_t\mid \mathcal{F}^0_t)$ of a representative state $X_t$ enters the drift and diffusion; see \cite{CarmonaDelarueMFGVol1,CarmonaDelarueMFGVol2,Lacker2018}.

\paragraph{Assumptions.}
We impose the following standard conditions (see, e.g., \cite{FournierGuillin2015,Lacker2018}).

\begin{enumerate}[label=(H\arabic*),ref=(H\arabic*)]
\item \textbf{(Lipschitz in state and measure.)}
The drift $b(t,x,\mu)$ and diffusion $\sigma(t,x,\mu)$ are Lipschitz in $(x,\mu)$ uniformly in $t\in[0,T]$, with $\mu$ in the space of probability measures on $\R^d$ metrized by the Wasserstein distance $W_2$.

\item \textbf{(Linear growth.)}
$|b(t,x,\mu)| + \|\sigma(t,x,\mu)\| \le C(1+|x| + W_2(\mu,\delta_0))$ for some $C\ge 0$.

\item \textbf{(Common noise integrability.)}
Common-noise coefficients are progressively measurable and square-integrable on $[0,T]$ a.s.

\item \textbf{(Initial condition.)}
The initial states $\{X^{i,N}_0\}_{i=1}^N$ are i.i.d.\ with law $\mu_0$ and satisfy $\E[|X^{i,N}_0|^2]<\infty$; $\mu_0$ has finite second moment.

\item \textbf{(Non-degeneracy / bounded coefficients.)}
The diffusion matrix $\sigma\sigma^\top$ is uniformly elliptic (or at least non-degenerate) and all coefficients are bounded when restricted to compacts in $\R^d\times \mathcal{P}_2(\R^d)$.
\end{enumerate}

Under (H1)--(H5), the limiting McKean--Vlasov equation admits a unique strong solution and the $N$-particle system is well-posed; see \cite{CarmonaDelarueMFGVol1,CarmonaDelarueMFGVol2}.

\begin{theorem}[Convergence of empirical measures]\label{thm:mfg_limit_conv}\label{thm:MFG_fixed_point}
Assume (H1)--(H5). Let $\{X^{i,N}\}_{i=1}^N$ solve the $N$-particle system and let $X$ solve the limiting McKean--Vlasov equation with the same common noise $W^0$ and i.i.d.\ copies $\{B^i\}$ of the idiosyncratic driver. Then for each $t\in[0,T]$,
\[
\bar\mu^N_t \xrightarrow[N\to\infty]{W_2} \mu_t \quad \text{in probability},
\]
where $\mu_t = \Law(X_t\mid \mathcal{F}^0_t)$ is the conditional law of the limit process. In particular, propagation of chaos holds: any fixed finite marginals converge in law to the product of the limit conditional law.
\end{theorem}

\begin{proof}[Proof sketch (black-box references)]
Existence and uniqueness for the limit equation follow from (H1)--(H2) and fixed-point arguments in the space of flows of measures; see \cite{CarmonaDelarueMFGVol1}. Convergence of $\bar\mu^N_t$ to $\mu_t$ in $W_2$ is a standard consequence of (H1)--(H5) and the conditional law of large numbers for the empirical measure; see \cite{FournierGuillin2015} for quantitative rates and \cite{Lacker2018} for the common-noise formulation. Propagation of chaos is then implied by the same estimates; see \cite{CarmonaDelarueMFGVol2}.
\end{proof}

\begin{proposition}[Stability of mean-field equilibrium]\label{prop:mfg_limit_stability}
Under (H1)--(H5), the mean-field equilibrium flow $\{\mu_t\}_{t\in[0,T]}$ is unique in the class of admissible flows with finite second moments. Moreover, $\sup_{t\in[0,T]}\E[W_2(\bar\mu^N_t,\mu_t)^2] \le C/N$ for some constant $C$ depending on $T$ and the data.
\end{proposition}

\begin{proof}[Proof sketch (black-box references)]
Uniqueness follows from the Lipschitz structure (H1) and Gronwall-type estimates in Wasserstein space; see \cite{CarmonaDelarueMFGVol1}. The $O(1/N)$ bound for $\E[W_2(\bar\mu^N_t,\mu_t)^2]$ is given in \cite{FournierGuillin2015} and \cite{HuangMalhameCaines2006} for the classical and MFG settings.
\end{proof}

\begin{remark}[LQG specialization]
In our LQG setting (see Section~IV), the coefficients are affine in state and measure, so (H1)--(H2) hold and the limit is explicit. The preceding black-box statements justify using the mean-field system as the large-$N$ approximation for the particle implementation in Section~VI.
\end{remark}

% LQG constructive solution: reduced state, fixed-point map, Picard algorithm.
% Compatible with ieee_main.tex (uses \E, \R, \CARA, etc.).

\subsection{Constructive LQG Solution via Picard Iteration}
\label{sec:lqg_constructive}

\paragraph{Benchmark (position impact; continuous time).}
This subsection is a continuous-time constructive solver for a \emph{position-impact} formulation (it collapses stock and flow so that $\bar u\equiv \bar x$).
It is included as a diagnostic/benchmark and is not used in the baseline discrete-time stock--flow simulations in Section~VI.

\paragraph{Endogenous vs exogenous.}
Mean demand $\bar x$, price $p$, and (from the equilibrium map) mispricing $\bar y = m - p$ and the coefficient $B_t$ (driven by $\bar x$) are \emph{endogenous}; in equilibrium, beliefs and the filter are consistent with the chosen information. Exogenous parameters include $\kappa$, $\lambda$, $\phi$, $\gamma$, and other primitives (e.g.\ $\sigma_v$, $\sigma_c$, $\sigma_\epsilon$, $k$). In this formulation, belief updating conditions only on the private signal $\xi$ and not on the price observation channel; hence the mean belief $m_t$ is not fed back by $\bar x$. Price is endogenous as an outcome of $\bar x$ but is not used as an observation in the filter; incorporating it would require joint filtering and a richer fixed-point formulation (see Section~IV).

\paragraph{Reduced state and dynamics.}
Define the (belief--price) mispricing state
\begin{equation}
y_t \;:=\; \hat v_t - p_t.
\label{eq:y_reduced}
\end{equation}
We work in the homogeneous representative-agent setting with overconfidence parameter $k$, so $\hat v_t$ is the subjective Kalman estimate of $v_t$, $\Sigma_t$ is the perceived posterior variance, and $R':=\sigma_c^2+\sigma_\epsilon^2/k$ is the perceived measurement-noise variance. Wealth and state dynamics in innovation form are
\begin{align}
dW_t
&= x_t\,\bigl(\kappa\,y_t + \lambda\,\bar x_t\bigr)\,dt + x_t\,\phi\,dI^p_t,
\label{eq:W_dyn_reduced}\\
dy_t
&= \bigl(\mu_v - \kappa\,y_t - \lambda\,\bar x_t\bigr)\,dt + \beta_t\,dI^\xi_t - \phi\,dI^p_t,
\label{eq:y_dyn_reduced}
\end{align}
where (under the agent's subjective filtering model in Appendix~A) $I^p$ and $I^\xi$ are independent Brownian motions (price and signal innovations). The (subjective) unhedgeable loading is
\begin{equation}
\beta_t \;=\; K_t\,\sqrt{\sigma_c^2+\sigma_\epsilon^2},\qquad
K_t \;=\; \frac{\Sigma_t}{R'},\qquad
R' \;=\; \sigma_c^2 + \frac{\sigma_\epsilon^2}{k},
\label{eq:beta_K_R}
\end{equation}
and the perceived posterior variance $\Sigma_t$ solves the Riccati equation
\begin{equation}
\dot\Sigma_t \;=\; \sigma_v^2 - \frac{\Sigma_t^2}{R'},\qquad \Sigma_0\ge 0.
\label{eq:Sigma_riccati}
\end{equation}

\paragraph{Fixed-point map and contraction.}
Theorem~\ref{thm:MFG_fixed_point} provides the framework: existence, uniqueness, and convergence of the mean-field limit are from the black-box theory (assumptions (H1)--(H5)); the reduced dynamics \eqref{eq:W_dyn_reduced}--\eqref{eq:y_dyn_reduced}, Riccati \eqref{eq:Sigma_riccati}, and mean-field update \eqref{eq:barx_update} below are derived in our LQG setting. Given a candidate mean-field demand process $\bar x(\cdot)$, one can (i) solve the backward Riccati for $A_t$ (independent of $\bar x$), (ii) solve the backward linear ODE for $B_t$ driven by $\bar x_t$, (iii) simulate forward the coupled $(v_t,m_t,p_t)$ and hence $\bar y_t = m_t - p_t$, and (iv) update $\bar x$ via the mean-field consistency relation. This defines a fixed-point map $\Phi$ on trajectories $\bar x(\cdot)$.

\paragraph{Fixed-point object.}
The fixed-point object is the map $\Phi$ on the space $\mathcal X := L^2_{\mathcal F^0}(\Omega\times[0,T])$ of candidate mean-field demand trajectories (adapted to the common filtration $\mathcal F^0$), equipped with the norm
\[
\|\bar x\|_2 := \Big(\E\int_0^T |\bar x_t|^2\,dt\Big)^{1/2}.
\]
Equilibrium corresponds to the unique fixed point $\bar x^\star = \Phi(\bar x^\star)$ when the contraction condition $q<1$ below holds; see Proposition~\ref{prop:lqg_constructive}.

\begin{proposition}[Constructive equilibrium under weak coupling / small horizon]
\label{prop:lqg_constructive}
Define the fixed-point map $\Phi$ on $\bar x(\cdot)$ as follows. Given $\bar x(\cdot)$:
\begin{enumerate}[label=(\roman*)]
\item Solve backward the Riccati equation for $A_t$ with terminal condition $A_T=0$ (as in the $A$-ODE of the exact \CARA/\LQG\ formulation).
\item Solve backward the linear ODE for $B_t$ driven by $\bar x_t$, with terminal condition $B_T=0$.
\item Solve forward the dynamics for $(v_t,m_t,p_t)$ given $\bar x_t$, and set $\bar y_t = m_t - p_t$.
\item Update via mean-field consistency:
\begin{equation}
\bar x_t \;=\; \frac{(\kappa + \phi^2 A_t)\,\bar y_t + \phi^2 B_t}{\gamma\,\phi^2 - \lambda}.
\label{eq:barx_update}
\end{equation}
\end{enumerate}
Define
\begin{equation}
q \;:=\; \frac{\lambda T}{\gamma\phi^2-\lambda}\,\bigl(2\kappa+\kappa^2 T\bigr).
\label{eq:q_contract}
\end{equation}
If $\gamma\phi^2>\lambda$ and $q<1$, then $\Phi$ is a contraction on the space $\mathcal X := L^2_{\mathcal F^0}(\Omega\times[0,T])$ and the fixed point exists and is unique; Picard iteration converges.
\end{proposition}

\begin{proof}
We make the contraction argument fully explicit.

\medskip\noindent
\textbf{Functional-analytic setup.}
Fix a filtered probability space $(\Omega,\mathcal F,\mathcal F^0,\Prob)$ supporting the common shocks
(e.g.\ $(W,U,Z)$ as in \eqref{eq:v_dyn_ct}--\eqref{eq:p_dyn_ct}) and the cross-sectional private noises.
Let $\mathcal F^0=(\mathcal F^0_t)_{t\in[0,T]}$ be the (common) market filtration containing the common shocks.
Define the Banach space
\[
\mathcal X := L^2_{\mathcal F^0}(\Omega\times[0,T]),\qquad
\|\bar x\|_2 := \Big(\E\int_0^T |\bar x_t|^2\,dt\Big)^{1/2}.
\]
(If one prefers to restrict to deterministic mean-field trajectories $\bar x:[0,T]\to\mathbb R$, the same proof applies
 verbatim with $\|\bar x\|_{L^2([0,T])}:=(\int_0^T|\bar x_t|^2dt)^{1/2}$.)

Assume $\gamma\phi^2>\lambda$, so the denominator in \eqref{eq:barx_update} is strictly positive.
Also assume $\beta$ is $\mathcal F^0$-progressively measurable and bounded on $[0,T]$ (this holds under the constant-gain
closure \eqref{eq:mean_field_closure}, or any bounded-gain filtering specification used in Sections III--IV).

\medskip\noindent
\textbf{Step 1: Riccati well-posedness and a uniform bound for $A$.}
In the homogeneous case, $A_t$ solves \eqref{eq:A_ode} with terminal condition $A_T=0$:
\[
-\dot A_t=\frac{\kappa^2}{\phi^2}-\beta_t^2 A_t^2,\qquad A_T=0.
\]
Set $\widetilde A_s := A_{T-s}$ for $s\in[0,T]$. Then
\[
\dot{\widetilde A}_s=\frac{\kappa^2}{\phi^2}-\beta_{T-s}^2\,\widetilde A_s^2,\qquad \widetilde A_0=0.
\]
\emph{Nonnegativity.} The vector field $f(s,a):=\kappa^2/\phi^2-\beta_{T-s}^2a^2$ is continuous in $s$ and locally
Lipschitz in $a$, hence solutions are unique. Moreover $f(s,0)=\kappa^2/\phi^2>0$ for all $s$, so starting from
$\widetilde A_0=0$ the solution becomes strictly positive immediately. If it ever became negative, let
$\tau:=\inf\{s>0:\widetilde A_s=0\}$ be the first return time to $0$ from above. Then $\dot{\widetilde A}_\tau=f(\tau,0)>0$,
so the trajectory points into the positive half-line at $\tau$, contradicting the definition of $\tau$. Thus
$\widetilde A_s\ge 0$ for all $s\in[0,T]$, hence $A_t\ge 0$ for all $t$.
\emph{Coarse bound.} Because the quadratic term is nonnegative,
\[
\dot{\widetilde A}_s \le \frac{\kappa^2}{\phi^2}
\quad\Longrightarrow\quad
0\le \widetilde A_s \le \frac{\kappa^2}{\phi^2}s,
\]
hence
\[
0\le A_t \le \frac{\kappa^2}{\phi^2}(T-t)\le \frac{\kappa^2}{\phi^2}T
\qquad\text{for all }t\in[0,T].
\]
In particular, $A\in L^\infty([0,T])$ and $\|A\|_\infty\le \kappa^2 T/\phi^2$.

\medskip\noindent
\textbf{Step 2: Precise assumption on mean belief and the $\bar y$-Lipschitz bound.}
Recall $\bar y_t := m_t-p_t$ in Proposition~III.7.

In the myopic formulation of Section~IV.B, belief updating (Appendix~A) is conditioned only on the private signal
$\xi_{i,t}$ (and not on the endogenous price observation channel). Under this modeling choice,
the mean belief process $m_t=\int_0^1\hat v_{i,t}\,di$ is independent of the mean-field demand input $\bar x$; i.e.\ for
two inputs $\bar x,\bar x'\in\mathcal X$ we have $m\equiv m'$.

(If one prefers not to hard-code this, it suffices to assume instead that $\bar x\mapsto m$ is Lipschitz in
$\|\cdot\|_2$ with constant $C_mT$; the bound below then holds with $\lambda$ replaced by $\lambda+C_m$.)

Now fix $\bar x,\bar x'\in\mathcal X$. Couple the two forward systems on the same probability space with the same common
shocks and identical initial conditions $(v_0,p_0)$. Subtracting the price equations \eqref{eq:p_dyn_ct} yields, pathwise,
\[
d(p_t-p_t') = -\kappa(p_t-p_t')\,dt + \lambda(\bar x_t-\bar x_t')\,dt,\qquad (p_0-p_0')=0,
\]
since the common noise term $\phi\,dZ_t$ cancels under this coupling.
Solving gives
\[
|p_t-p_t'|
\le \lambda\int_0^t e^{-\kappa(t-s)}|\bar x_s-\bar x_s'|\,ds
\le \lambda\int_0^t|\bar x_s-\bar x_s'|\,ds
\le \lambda\sqrt{t}\Big(\int_0^t|\bar x_s-\bar x_s'|^2\,ds\Big)^{1/2}.
\]
Squaring and integrating in $t$, then using Fubini, yields
\[
\int_0^T |p_t-p_t'|^2\,dt
\le \lambda^2\int_0^T t\int_0^t|\bar x_s-\bar x_s'|^2\,ds\,dt
\le \lambda^2T^2\int_0^T|\bar x_s-\bar x_s'|^2\,ds.
\]
Taking expectations gives $\|p-p'\|_2\le \lambda T\|\bar x-\bar x'\|_2$.
Because $m\equiv m'$, we have $\bar y-\bar y'=-(p-p')$, hence
\[
\|\bar y-\bar y'\|_2 \le \lambda T\|\bar x-\bar x'\|_2.
\]

\medskip\noindent
\textbf{Step 3: $B$-Lipschitz bound in $\bar x$.}
Given $A$ (which does not depend on $\bar x$) and $\beta$ (independent of $\bar x$ under the same information closure as
above), $B_t$ solves \eqref{eq:B_ode} with terminal condition $B_T=0$:
\[
-\dot B_t=\frac{\kappa\lambda}{\phi^2}\bar x_t+\mu_v A_t-\beta_t^2A_tB_t,\qquad B_T=0.
\]
(If $\beta$ is allowed to depend on $\bar x$, an additional Lipschitz term appears; the same contraction logic goes
through with a larger constant.)

Let $B,B'$ be the solutions corresponding to $\bar x,\bar x'$ and set $\Delta B:=B-B'$.
Subtracting the ODEs gives
\[
-\dot{\Delta B}_t=\frac{\kappa\lambda}{\phi^2}(\bar x_t-\bar x_t')-\beta_t^2A_t\,\Delta B_t,\qquad \Delta B_T=0.
\]
By integrating factor,
\[
\Delta B_t=\frac{\kappa\lambda}{\phi^2}\int_t^T
\exp\!\Big(-\int_t^s \beta_u^2A_u\,du\Big)\,(\bar x_s-\bar x_s')\,ds.
\]
Since $A_u\ge 0$ and $\beta_u^2\ge 0$, the exponential factor is at most $1$, yielding
\[
|B_t-B_t'|
\le \frac{\kappa\lambda}{\phi^2}\int_t^T|\bar x_s-\bar x_s'|\,ds
\le \frac{\kappa\lambda}{\phi^2}\sqrt{T-t}\Big(\int_t^T|\bar x_s-\bar x_s'|^2\,ds\Big)^{1/2}.
\]
Squaring and integrating in $t$, then using Fubini, yields $\|B-B'\|_2 \le (\kappa\lambda/\phi^2)\,T\,\|\bar x-\bar x'\|_2$.

\medskip\noindent
\textbf{Step 4: Contraction estimate for $\Phi$ and a fully checkable sufficient condition.}
Recall \eqref{eq:barx_update}:
\[
\Phi(\bar x)_t
=
\frac{(\kappa+\phi^2A_t)\bar y_t+\phi^2 B_t}{\gamma\phi^2-\lambda}.
\]
Using $\gamma\phi^2-\lambda>0$, boundedness of $A$, and the triangle inequality,
\[
\|\Phi(\bar x)-\Phi(\bar x')\|_2
\le
\frac{(\kappa+\phi^2\|A\|_\infty)\|\bar y-\bar y'\|_2+\phi^2\|B-B'\|_2}{\gamma\phi^2-\lambda}.
\]
Substituting the bounds from Steps 2--3 gives
\[
\|\Phi(\bar x)-\Phi(\bar x')\|_2
\le
\frac{\lambda T}{\gamma\phi^2-\lambda}\Big(2\kappa+\phi^2\|A\|_\infty\Big)\,\|\bar x-\bar x'\|_2.
\]
Finally, using $\|A\|_\infty\le \kappa^2T/\phi^2$ from Step 1 yields the parameter-only bound
\[
\|\Phi(\bar x)-\Phi(\bar x')\|_2
\le
q\,\|\bar x-\bar x'\|_2,
\qquad
q:=\frac{\lambda T}{\gamma\phi^2-\lambda}\Big(2\kappa+\kappa^2T\Big),
\]
which is fully checkable from primitives.

\medskip\noindent
\textbf{Step 5: Banach fixed point.}
If $q<1$, then $\Phi$ is a contraction on $(\mathcal X,\|\cdot\|_2)$, hence admits a unique fixed point
$\bar x^\star\in\mathcal X$. Moreover, the Picard iteration $\bar x^{n+1}=\Phi(\bar x^n)$ converges geometrically to
$\bar x^\star$ in $\|\cdot\|_2$.

\medskip\noindent
\textbf{Remark (relation to the simpler myopic scaling).}
In the myopic closure where hedging terms are suppressed ($A\equiv 0$, $B\equiv 0$), the Lipschitz estimate tightens to
a pure $O(T)$ scaling proportional to $\kappa\lambda/(\gamma\phi^2-\lambda)$, matching the heuristic
small-horizon/weak-coupling condition discussed around \eqref{eq:q_contract}.
\end{proof}

\paragraph{Picard iteration.}
A conservative Lipschitz bound is $q$ in \eqref{eq:q_contract}; for $q<1$, Banach yields uniqueness.
Algorithm~\ref{alg:picard} summarizes the procedure. We initialize a candidate $\bar x^0(\cdot)$ (e.g., identically zero or the myopic closure $\bar x_t^0 = \kappa\bar y_t^0/(\gamma\phi^2-\lambda)$ from a forward run with $\bar x\equiv 0$), then iterate: compute $\Sigma,K,\beta$ from \eqref{eq:Sigma_riccati}--\eqref{eq:beta_K_R}, solve $A$ backward, solve $B$ backward using the current $\bar x^n$, forward-simulate $(v,m,p)$ and $\bar y^n$, update $\bar x^{n+1}$ via \eqref{eq:barx_update}, and stop when the residual (e.g., $\|\bar x^{n+1}-\bar x^n\|$ in $L^2$ or sup-norm over the time grid) is below a tolerance.

\begin{algorithm}[t]
\caption{Picard iteration for LQG mean-field equilibrium}
\label{alg:picard}
\begin{algorithmic}[1]
\STATE \textbf{Input:} Horizon $T$, parameters $(\kappa,\lambda,\gamma,\phi,\mu_v,\sigma_v,\sigma_c,\sigma_\epsilon,k)$, initial conditions $(v_0,p_0,\hat v_0,\Sigma_0)$, tolerance $\epsilon$, max iterations $N_{\max}$.
\STATE \textbf{Initialize:} Set $\bar x^0_t \gets 0$ (or myopic seed) for $t\in[0,T]$. Set $n\gets 0$.
\REPEAT
\STATE Compute $\Sigma_t$, $K_t$, $\beta_t$ on $[0,T]$ from \eqref{eq:Sigma_riccati} and \eqref{eq:beta_K_R}.
\STATE Solve backward for $A_t$ (Riccati, $A_T=0$).
\STATE Solve backward for $B_t$ given $\bar x^n_t$ (linear ODE, $B_T=0$).
\STATE Forward-simulate $(v_t,m_t,p_t)$ and $\bar y^n_t = m_t - p_t$ using $\bar x^n_t$ in the dynamics.
\STATE Update $\bar x^{n+1}_t \gets \frac{(\kappa+\phi^2 A_t)\bar y^n_t + \phi^2 B_t}{\gamma\phi^2-\lambda}$ for $t\in[0,T]$.
\STATE Compute residual $r_n \gets \|\bar x^{n+1}-\bar x^n\|$ (e.g., $L^2$ or sup over grid).
\STATE $n \gets n+1$.
\UNTIL{$r_{n-1}<\epsilon$ or $n\ge N_{\max}$}
\STATE \textbf{Output:} Equilibrium mean-field demand $\bar x^* \approx \bar x^n$ and associated $(A,B,\bar y)$.
\end{algorithmic}
\end{algorithm}

\paragraph{Implementation and reproducibility.}
The Picard iteration in Algorithm~\ref{alg:picard} is implemented in a MATLAB script; the path and run instructions will be specified in the replication package. All reported equilibrium paths and statistics use this routine with a fixed tolerance and time discretization.

\subsection{Reproducibility and numerical verification}
The equilibrium is computed via Algorithm 1 implemented in \ifjeicblind the replication package (link omitted for blind review)\else \path{code/matlab/mfg_lqg_common_noise_demo.m}\fi.
Diagnostics include the fixed-point residual along Picard iterates, a time-step refinement test under decreasing $dt$, and an $N$-convergence test comparing particle moments to the mean-field limit.\footnote{Outputs and plots can be regenerated by running the script.}

\paragraph{Remark (discrete-time certification).}
At the discrete-time level, uniqueness of the numerical fixed point can be certified via a Banach fixed-point argument on trajectories endowed with a sup-norm metric.
We provide a machine-checkable backbone for this contraction argument \ifjeicblind (link omitted for blind review)\else in \texttt{lean/MFGContraction.lean}\fi; it does not attempt to formalize the full stochastic propagation-of-chaos theorem.


\section{Main Results}
\label{sec:results}
% Section 4: Main Results

\subsection{Static Benchmark: Price Sensitivity and Overconfidence}

The closed-form static price \eqref{eq:static_price} yields explicit comparative statics with respect to overconfidence.

\begin{proposition}[Static price response increases with overconfidence]
\label{prop:static_comparative}
Let $\tau_{v,0} := 1/\sigma_{v,0}^2$ and $\tau'_{\epsilon,0} := k/\sigma_{\epsilon,0}^2$. In the static benchmark, the coefficient on $v$ in \eqref{eq:static_price} is
\begin{equation}
a(k) = \frac{\lambda \tau'_{\epsilon,0}}{\gamma + \lambda(\tau_{v,0}+\tau'_{\epsilon,0})},
\end{equation}
which is increasing in $k$ and satisfies $0<a(k)<1$.
\end{proposition}

\subsection{Myopic Trading Rule and Mean-Field Consistency}

We next derive a checkable \emph{myopic} inventory rule for the discrete-time timing in Subsubsection~\ref{subsubsec:timing_stock_flow}. Fix a decision interval $\Delta t$ and treat the mean trading rate $\bar u_t$ and the public price $p_t$ as given. For tractability, belief updating (Appendix~A) conditions only on the private signal $\xi_i$; incorporating endogenous prices as an additional observation channel would require joint filtering and a richer fixed-point formulation. Conditional on the investor's information $\mathcal F^i_t$, the one-step price increment $\Delta p_{t+\Delta t}$ is Gaussian with conditional mean $(\kappa(\hat v_{i,t}-p_t)+\lambda \bar u_t)\Delta t$ and conditional variance $\sigma^2_{p,i,t}\Delta t$.

\begin{proposition}[Myopic \CARA inventory choice with holding cost and flow impact]
\label{prop:myopic_cara}
Fix a decision interval $\Delta t$. Suppose conditional on $\mathcal F^i_t$ the price increment is Gaussian with
$\E_i[\Delta p_{t+\Delta t}]=(\kappa(\hat v_{i,t}-p_t)+\lambda\bar u_t)\Delta t$ and
$\Var_i(\Delta p_{t+\Delta t})=\sigma^2_{p,i,t}\Delta t$.
With one-step wealth $\Delta W_{i,t}=x_{i,t}\Delta p_{t+\Delta t}-(\chi/2)x_{i,t}^2\Delta t$ and \CARA utility, the unique maximizer is
\begin{equation}
x^\star_{i,t} = \frac{\kappa(\hat v_{i,t}-p_t)+\lambda\bar u_t}{\chi+\gamma\,\sigma^2_{p,i,t}}.
\label{eq:myopic_policy}
\end{equation}
\end{proposition}

Define the corresponding \emph{mispricing gain} (shares per unit price) as the coefficient on $(\hat v_{i,t}-p_t)$ in \eqref{eq:myopic_policy}. Under the Euler increment \eqref{eq:anchored_impact} over $\Delta t$ and conditioning on $\mathcal F^i_t$, $\sigma^2_{p,i,t}=\phi^2+\kappa^2\Sigma_{i,t}\Delta t$, so
\begin{equation}
\Theta_t(\Sigma_{i,t})
:=
\frac{\kappa}{\chi+\gamma\,\sigma^2_{p,i,t}}
=
\frac{\kappa}{\chi+\gamma\bigl(\phi^2+\kappa^2\Sigma_{i,t}\Delta t\bigr)}.
\label{eq:theta_def}
\end{equation}
Moreover,
\[
\partial_\Sigma \Theta_t(\Sigma)
=
-\frac{\gamma \kappa^3 \Delta t}{\bigl(\chi+\gamma(\phi^2+\kappa^2\Sigma\Delta t)\bigr)^2}\le 0,
\]
so $\Theta_t$ is weakly decreasing in $\Sigma$.

\paragraph{Position-impact benchmark (reduced form; not used in baseline simulations).}
If one collapses stock and flow (so $\bar u_t\equiv \bar x_t$) and sets $\chi=0$ and $\sigma^2_{p,i,t}\equiv\phi^2$, then averaging \eqref{eq:myopic_policy} yields the familiar homogeneous closure (assuming $\gamma\phi^2>\lambda$):
\begin{equation}
\bar x_t = \frac{\kappa}{\gamma\phi^2-\lambda}\,(m_t-p_t),
\label{eq:mean_field_closure}
\end{equation}
where $m_t := \int_0^1 \hat v_{i,t}\,di$ is the mean belief. This reduced form is retained only as a transparent diagnostic benchmark; the baseline simulations solve the discrete-time stock--flow fixed point with inventory carryover and uniqueness condition $0\le B_t<\Delta t$ (Proposition~\ref{prop:step_fixed_point} and Theorem~\ref{thm:contraction}).

\subsection{Exact continuous-time CARA/LQG control with filtering and mean-field coupling}
\label{subsec:ct_cara_lqg}

This subsection replaces the infinitesimal-horizon (myopic) CARA rule with the
finite-horizon continuous-time exponential-utility optimum under partial information.
We keep the same fundamental/signal/price dynamics \eqref{eq:v_dyn_ct}--\eqref{eq:p_dyn_ct}
and the same Kalman--Bucy filtering structure for subjective beliefs.

\paragraph{Dynamics and information.}
Let the fundamental and price evolve as
\begin{align}
dv_t &= \mu_v\,dt + \sigma_v\,dW_t, \label{eq:v_dyn_ct}\\
d\xi_{i,t} &= v_t\,dt + \sigma_c\,dU_t + \sigma_\epsilon\,dB_{i,t}, \label{eq:xi_dyn_ct}\\
dp_t &= \kappa\,(v_t - p_t)\,dt + \lambda\,\bar u_t\,dt + \phi\,dZ_t, \label{eq:p_dyn_ct}
\end{align}
where $\bar u_t := \int_0^1 u_{j,t}\,dj$ is the mean-field trading rate (taken as exogenous by an infinitesimal agent).
Agent $i$ observes $(p_s,\xi_{i,s})_{s\le t}$ and forms a subjective filtered estimate $\hat v_{i,t}$.

\paragraph{Subjective Kalman filter (with overconfidence).}
As in Appendix~A, agent $i$ runs a (possibly misspecified) Kalman--Bucy filter
\begin{equation}
\begin{aligned}
d\hat v_{i,t}
&= \mu_v\,dt + K_t^{(i)}\bigl(d\xi_{i,t}-\hat v_{i,t}\,dt\bigr),\\
K_t^{(i)} &:= \frac{\Sigma_t^{(i)}}{R'_i},\qquad
R'_i := \sigma_c^2+\frac{\sigma_\epsilon^2}{k_i},
\end{aligned}
\label{eq:kb_filter_ct}
\end{equation}
with Riccati equation
\begin{equation}
\dot\Sigma_t^{(i)} \;=\; \sigma_v^2 - \frac{(\Sigma_t^{(i)})^2}{R'_i},
\qquad \Sigma_0^{(i)}\ge 0.
\label{eq:riccati_ct}
\end{equation}
(When convenient, one may use the constant-gain approximation $K_t^{(i)}\approx K_i^\star
= \sigma_v/\sqrt{R'_i}$.)

\paragraph{Innovation form and reduced state.}
Define the (belief-based) mispricing state
\begin{equation}
y_{i,t} := \hat v_{i,t}-p_t.
\label{eq:y_def}
\end{equation}
Under the agent's filtration, write the observed price in innovations form
\begin{equation}
dp_t \;=\; \bigl[\kappa\,y_{i,t}+\lambda\,\bar u_t\bigr]dt \;+\; \phi\,dI^p_{i,t},
\label{eq:p_innov}
\end{equation}
where $I^p_{i}$ is a Brownian motion (price innovation).
Likewise, the (normalized) signal innovation process is defined by
\[
dI^\xi_{i,t} \;:=\; \frac{d\xi_{i,t}-\hat v_{i,t}\,dt}{\sqrt{\sigma_c^2+\sigma_\epsilon^2}}.
\]
Under the agent filtration $\mathcal F^i_t$, $I^\xi_i$ is a Brownian motion and is independent of the price innovation $I^p_i$.

Consequently, the controlled wealth and the state satisfy
\begin{align}
dW_{i,t} &= x_{i,t}\,dp_t \;=\; x_{i,t}\bigl[\kappa\,y_{i,t}+\lambda\,\bar u_t\bigr]dt
          + x_{i,t}\phi\,dI^p_{i,t}, \label{eq:W_dyn}\\
dy_{i,t} &= \bigl[\mu_v-\kappa\,y_{i,t}-\lambda\,\bar u_t\bigr]dt
          + \beta^{(i)}_t\,dI^\xi_{i,t} \;-\; \phi\,dI^p_{i,t},
\label{eq:y_dyn}
\end{align}
where $I^\xi_i$ is an (independent) Brownian motion and the (subjective) unhedgeable
belief-noise loading is
\begin{equation}
\beta^{(i)}_t \;:=\; K_t^{(i)}\,\sqrt{\sigma_c^2+\sigma_\epsilon^2}.
\label{eq:beta_def}
\end{equation}
Note that $K_t^{(i)}$ depends on the perceived measurement-noise variance $R'_i(k_i)$, whereas the innovation volatility factor $\sqrt{\sigma_c^2+\sigma_\epsilon^2}$ reflects the true signal-noise variance in $d\xi_{i,t}$.

\medskip
Equivalently, since $d\xi_{i,t}-\hat v_{i,t}\,dt=\sqrt{\sigma_c^2+\sigma_\epsilon^2}\,dI^\xi_{i,t}$, the Kalman--Bucy update \eqref{eq:kb_filter_ct} can be written in innovation form as
\[
d\hat v_{i,t}=\mu_v\,dt+\beta^{(i)}_t\,dI^\xi_{i,t}.
\]
This is the diffusion term that enters the reduced state dynamics \eqref{eq:y_dyn}.

(Under constant gain, $\beta^{(i)}_t\approx \beta_i^\star = K_i^\star\sqrt{\sigma_c^2+\sigma_\epsilon^2}$
is constant.)

\paragraph{Diffusions entering the HJB.}
In the innovation representation \eqref{eq:W_dyn}--\eqref{eq:y_dyn}, the controlled state is $(W_{i,t},y_{i,t})$ and is driven by two independent Brownian motions: the \emph{price innovation} $I^p_{i}$ with volatility $\phi$ and the \emph{filter innovation} $I^\xi_{i}$ with volatility $\beta^{(i)}_t$.
Consequently, the HJB/It\^o generator contains the variance terms $\frac12 x_{i,t}^2\phi^2 V_{ww}$ and $\frac12\bigl((\beta^{(i)}_t)^2+\phi^2\bigr)V_{yy}$.
Moreover, because the same $dI^p_{i,t}$ appears in both $dW_{i,t}$ and $dy_{i,t}$, we have the cross-variation
$d\langle W_i,y_i\rangle_t = -x_{i,t}\phi^2\,dt$,
which produces the mixed derivative term $-x_{i,t}\phi^2 V_{wy}$ in the HJB and is exactly the source of the intertemporal hedging correction in \eqref{eq:opt_control_lqg}.

\paragraph{Control objective.}
Agent $i$ maximizes exponential utility of terminal wealth:
\begin{equation}
\sup_{x_i\in\mathcal A}\;
\mathbb E\Bigl[-\exp\bigl(-\gamma_i W_{i,T}\bigr)\,\big|\,\mathcal F^i_t\Bigr],
\label{eq:exp_obj}
\end{equation}
over admissible progressively measurable controls $\mathcal A$ with
$\mathbb E\!\int_0^T x_{i,t}^2\,dt < \infty$.

\paragraph{Exact CARA/LQG solution (finite horizon).}
\begin{proposition}[Exponential-quadratic value and optimal feedback]
\label{prop:ct_cara_lqg}
Fix a mean-field trading-rate process $(\bar u_t)_{t\in[0,T]}$.
The value function admits the exponential-quadratic form
\begin{equation}
V_i(t,w,y) \;=\; -\exp\!\left(-\gamma_i w - \frac12 A^{(i)}_t y^2 - B^{(i)}_t y - C^{(i)}_t\right),
\label{eq:V_form}
\end{equation}
where $(A^{(i)}_t,B^{(i)}_t,C^{(i)}_t)$ solve the backward ODE system with terminal conditions
$A^{(i)}_T=B^{(i)}_T=C^{(i)}_T=0$:
\begin{align}
-\dot A^{(i)}_t
&= \frac{\kappa^2}{\phi^2} - \bigl(\beta^{(i)}_t\bigr)^2\bigl(A^{(i)}_t\bigr)^2,
\label{eq:A_ode}\\
-\dot B^{(i)}_t
 &= \frac{\kappa\lambda}{\phi^2}\,\bar u_t \;+\; \mu_v\,A^{(i)}_t
    \;-\; \bigl(\beta^{(i)}_t\bigr)^2 A^{(i)}_t B^{(i)}_t,
\label{eq:B_ode}\\
-\dot C^{(i)}_t
&= \frac12\Bigl(\bigl((\beta^{(i)}_t)^2+\phi^2\bigr)A^{(i)}_t
      - (\beta^{(i)}_t)^2\bigl(B^{(i)}_t\bigr)^2\Bigr)\\
&\quad + \mu_v B^{(i)}_t + \frac{\lambda^2}{2\phi^2}\bar u_t^{\,2}.
\label{eq:C_ode}
\end{align}
The optimal control is the affine feedback
\begin{equation}
x_{i,t}^\star
\;=\; \frac{\kappa\,y_{i,t}+\lambda\,\bar u_t}{\gamma_i\phi^2}
      \;+\; \frac{A^{(i)}_t y_{i,t}+B^{(i)}_t}{\gamma_i}.
\label{eq:opt_control_lqg}
\end{equation}
The first term is the myopic (instantaneous) demand.
The second term is an intertemporal hedging correction induced by the correlation between the state $y_{i,t}$
and traded return noise in \eqref{eq:y_dyn}.
\end{proposition}

\paragraph{Constant-gain closed form (tight approximation).}
Under the constant-gain approximation $K_t^{(i)}\approx K_i^\star$ so that
$\beta^{(i)}_t\approx \beta_i^\star$ is constant, \eqref{eq:A_ode} has the explicit solution
\begin{equation}
\begin{aligned}
A^{(i)}_t
&\approx
\frac{\kappa}{\phi\,\beta_i^\star}\,
\tanh\!\left(\frac{\kappa\,\beta_i^\star}{\phi}(T-t)\right),\\
\beta_i^\star &:= K_i^\star\sqrt{\sigma_c^2+\sigma_\epsilon^2},\\
K_i^\star &:= \frac{\sigma_v}{\sqrt{\sigma_c^2+\sigma_\epsilon^2/k_i}}.
\end{aligned}
\label{eq:A_closed_form}
\end{equation}
The remaining $(B^{(i)}_t,C^{(i)}_t)$ follow from linear ODEs once $\bar u_t$ is specified.

\paragraph{Mean-field consistency.}
With flow-based impact, the endogenous mean-field object is the trading rate $\bar u_t$ rather than the inventory $\bar x_t$.
Closing the loop therefore requires an explicit link between inventory adjustments and trading rates; Subsection~\ref{subsubsec:timing_stock_flow} fixes the discrete-time timing used in the baseline particle implementation, and Subsection~\ref{subsec:ct_flow_control_extension} sketches a continuous-time flow-control extension.

\subsection{Continuous-time CARA/LQG with inventory dynamics and quadratic trading costs (flow-control extension)}
\label{subsec:ct_flow_control_extension}
This subsection provides a continuous-time benchmark consistent with the stock--flow separation (inventory $x$ versus trading rate $u$). It is included only as an extension and is not used in the baseline discrete-time simulations.

Let $u_{i,t}$ denote the agent's trading rate and let inventory evolve as
\begin{equation}
dx_{i,t}=u_{i,t}\,dt.
\end{equation}
Under the innovation representation (Proposition~\ref{prop:ct_cara_lqg}), price evolves with flow-based impact,
\begin{equation}
dp_t = (\kappa y_{i,t}+\lambda \bar u_t)\,dt+\phi\,dI^p_{i,t},
\end{equation}
and we augment wealth with quadratic trading and holding costs,
\begin{equation}
dW_{i,t}=x_{i,t}\,dp_t-\eta\,u_{i,t}^2\,dt-\frac{\chi}{2}x_{i,t}^2\,dt,
\qquad \eta>0,\;\chi>0.
\end{equation}
Consider the exponential-utility value function with state $(w,x,y)$,
\[
V(t,w,x,y):=\sup_{(u_s)_{s\in[t,T]}} \E\bigl[-\exp(-\gamma W_{i,T})\mid (W_{i,t},x_{i,t},y_{i,t})=(w,x,y)\bigr],
\]
taking the mean-field path $\bar u_t$ as given. With the CARA ansatz
\[
V(t,w,x,y)=-\exp\bigl(-\gamma(w+f(t,x,y))\bigr),
\]
the HJB takes the standard risk-sensitive LQG form, and the control enters only through
\[
\sup_{u}\,\bigl\{u\,\partial_x f(t,x,y)-\eta\,u^2\bigr\},
\]
so the optimal trading rate is
\begin{equation}
u_{i,t}^\star=\frac{\partial_x f(t,x_{i,t},y_{i,t})}{2\eta}.
\end{equation}
Under a quadratic specification for $f(t,x,y)$, $\partial_x f$ is affine in $(x,y)$ and hence $u_{i,t}^\star$ is linear in $(x,y)$. Substituting the quadratic ansatz into the HJB yields a coupled system of Riccati ODEs for the quadratic coefficients and linear ODEs for the affine terms, with forcing that depends on the (given) mean-field flow $\bar u_t$; this is the usual LQG structure.

\subsection{Level Amplification vs.\ Overconfidence Effects}
\label{subsec:amplification}

This subsection distinguishes two notions: (i) \emph{level amplification}---large mispricing/volatility levels driven by weak anchoring and/or high noise trading; and (ii) \emph{overconfidence (treatment) effects}---incremental differences when increasing $k$ from $1$ to $3$. The following proposition gives an analytical mechanism for why $k$ predominantly affects disagreement rather than mispricing.

\begin{proposition}[Mechanism: mean-belief closure and why $k$ shifts dispersion more than price]
\label{prop:mechanism_k}
Within the myopic mean-field equilibrium (per-step stock--flow closure in Proposition~\ref{prop:step_fixed_point}):
\begin{enumerate}[label=(\roman*)]
\item \textbf{Closure depends only on the mean belief.}
The per-step fixed point for $(\bar x_t,\bar u_t)$ depends on the belief distribution only through cross-sectional averages of $\hat v_{i,t}-p_t$ (equivalently, the mean belief $m_t:=\int_0^1 \hat v_{i,t}\,di$) and the predetermined pre-trade mean inventory $\bar x_{t^-}$. In particular, cross-sectional variance (disagreement) does not enter the mean-field fixed point except through the risk denominators $D_{i,t}$ (and in the homogeneous benchmark $D_{i,t}$ is common across agents).
(In the particle implementation, the one-step myopic rule uses the effective variance
$\mathrm{Var}^{\mathrm{eff}}_{i,t}=\phi^2+\kappa^2\Sigma_{i,t}\Delta t$ (Proposition~\ref{prop:step_fixed_point}); in the homogeneous benchmark $\Sigma_{i,t}$ is common across agents, so this term only rescales $\bar x_t$ and does not introduce dependence on disagreement.)
\item \textbf{Decomposition of mispricing.}
For additional closed-form intuition, under the reduced-form position-impact benchmark closure \eqref{eq:mean_field_closure} (which collapses stock and flow), let $y_t := p_t - v_t$ (mispricing) and $e_t := m_t - v_t$ (mean-belief tracking error). Under the constant-gain approximation, the pair $(y_t,e_t)$ satisfies the linear system \eqref{eq:2d_ou}: mispricing dynamics are driven by (a) \emph{microstructure regime} $(\kappa,\lambda,\phi)$ through the coefficient $c := \lambda\kappa/(\gamma\phi^2-\lambda)$ and the mean-reversion rate $\kappa+c$; and (b) the \emph{mean-belief term} $e_t$, which is the only channel through which the belief distribution affects price. Thus $\E[|p-v|]$ and stationary mispricing variance are determined by $(\kappa,\lambda,\phi)$ and the stationary law of $e_t$.
\item \textbf{Why $k$ predominantly affects disagreement.}
Under the information structure of Appendix~A, each agent's belief has a common component (loading $K^*(k)\sigma_c$ on $dU_t$) and an idiosyncratic component (loading $K^*(k)\sigma_\epsilon$ on $dB_{i,t}$). The cross-sectional mean $m_t$ aggregates beliefs, so the idiosyncratic components average out and $m_t$ is driven only by the common noise and the fundamental. By Proposition~\ref{prop:stationary_disagreement}, disagreement $D_\infty(k)=\Var(\hat v_{i,t}-m_t\mid \mathcal F^0)$ is \emph{strictly increasing} in $k$. By Proposition~\ref{prop:tracking_error}, the mean-belief error $e_t=m_t-v_t$ has stationary variance $\Var(e_\infty)$ that depends on $k$ only through $K^*(k)$ and is \emph{strictly decreasing} in $k$. Hence incremental increases in $k$ raise cross-sectional variance (disagreement) and trading intensity by construction, while the mean belief---and therefore equilibrium demand and mispricing---respond only weakly because they depend on $e_t$, whose variance is bounded and not uniformly amplified in $k$.
\end{enumerate}
\end{proposition}

\paragraph{Structural extensions (potential mechanisms).}
The model can be extended with structural mechanisms that could amplify mispricing persistence or volatility:
(i) \emph{State-dependent anchoring}: $\kappa(t) = \kappa_{\text{base}} \exp(-\alpha_\kappa |p_t - v_t|)$ weakens arbitrage pressure under extreme mispricing (limits to arbitrage);
(ii) \emph{Endogenous fundamental feedback}: $dv_t = \mu_v\,dt + \sigma_v\,dW_t + \beta_{\text{feedback}} (p_t - v_t)\,dt$ creates reflexive price--fundamental loops;
(iii) \emph{Overconfidence-dependent impact}: $\lambda_{\text{eff}} = \lambda_{\text{base}} (1 + \beta_\lambda (\bar k_t - 1))$ makes impact increase with aggregate overconfidence.

\begin{proposition}[Per-step fixed point in the particle implementation]
\label{prop:step_fixed_point}
In the discrete-time particle implementation, we use an effective variance
\begin{equation}
\mathrm{Var}^{\mathrm{eff}}_{i,t} := \phi^2 + \kappa^2\Sigma_{i,t}\,\Delta t,
\end{equation}
i.e.\ we instantiate the variance rate in Proposition~\ref{prop:myopic_cara} as $\sigma^2_{p,i,t}=\mathrm{Var}^{\mathrm{eff}}_{i,t}$,
and define the effective risk/holding-cost term
\begin{equation}
D_{i,t}:=\chi+\gamma\,\mathrm{Var}^{\mathrm{eff}}_{i,t}.
\end{equation}
The myopic rule takes the linear form
\begin{equation}
x_{i,t} = \frac{\alpha_0}{D_{i,t}}\Bigl(\kappa(\hat v_{i,t}-p_t) + \lambda \bar u_t\Bigr).
\label{eq:impl_policy}
\end{equation}
This form follows by applying Proposition~\ref{prop:myopic_cara} to the flow-based Euler increment
$\Delta p_{t+1}:=p_{t+1}-p_t=\kappa(v_t-p_t)\Delta t+\lambda \bar u_t\Delta t+\phi\sqrt{\Delta t}\,\varepsilon^\eta_{t+1}$:
conditional on $\mathcal F^i_t$, the drift depends on $\hat v_{i,t}$ and the variance satisfies
$\Var_{i,t}(\Delta p_{t+1})=\phi^2\Delta t+\kappa^2\Sigma_{i,t}(\Delta t)^2=\mathrm{Var}^{\mathrm{eff}}_{i,t}\,\Delta t$.
Hence the optimal one-step inventory depends on the effective variance rate $\mathrm{Var}^{\mathrm{eff}}_{i,t}=\phi^2+\kappa^2\Sigma_{i,t}\Delta t$.
The scalar $\alpha_0>0$ is an optional position-size (risk-budget) scaling (set to $\alpha_0=1$ in our baseline) and can be absorbed into a rescaling of $\gamma$.
Define
\[
A_t := \frac1N\sum_{i=1}^N \frac{\alpha_0\kappa(\hat v_{i,t}-p_t)}{D_{i,t}},
\qquad
B_t := \frac1N\sum_{i=1}^N \frac{\alpha_0\lambda}{D_{i,t}}.
\]
Let $\bar x_{t^-}:=\frac1N\sum_{i=1}^N x_{i,t^-}$ denote the pre-trade mean inventory and note that the mean trading rate satisfies
$\bar u_t=(\bar x_t-\bar x_{t^-})/\Delta t$.
If $0\le B_t<\Delta t$, the scalar fixed point exists and is unique and is given by
\begin{equation}
\bar x_t = \frac{A_t - (B_t/\Delta t)\bar x_{t^-}}{1 - B_t/\Delta t},
\qquad
\bar u_t=\frac{\bar x_t-\bar x_{t^-}}{\Delta t}.
\end{equation}
\end{proposition}

\begin{theorem}[Contraction and uniqueness of the per-step mean-field closure]
\label{thm:contraction}
Fix time $t$ and treat the belief states $(\hat v_{i,t},\Sigma_{i,t})$, the public price $p_t$, and the pre-trade mean inventory $\bar x_{t^-}$ as given. Define the affine mean-field map
\begin{equation}
\Phi_t(\bar x) := A_t + \frac{B_t}{\Delta t}\bar x - \frac{B_t}{\Delta t}\bar x_{t^-},
\end{equation}
where $A_t,B_t$ are as in Proposition~\ref{prop:step_fixed_point}. If $0\le B_t<\Delta t$, then $\Phi_t$ is a contraction on $\mathbb R$ with constant $B_t/\Delta t$ and has a unique fixed point, given in closed form in Proposition~\ref{prop:step_fixed_point}.
Moreover, since $\mathrm{Var}^{\mathrm{eff}}_{i,t}\ge \phi^2$ and $D_{i,t}\ge \chi+\gamma\phi^2$, a simple sufficient condition ensuring $B_t<\Delta t$ uniformly is
\begin{equation}
\frac{\alpha_0\lambda}{\Delta t\,(\chi+\gamma\phi^2)}<1.
\label{eq:global_contraction}
\end{equation}
In the homogeneous case $D_{i,t}\equiv D_t$, we have $B_t=\alpha_0\lambda/D_t$, so $0\le B_t<\Delta t$ is equivalent to
$D_t>\alpha_0\lambda/\Delta t$, i.e.\ $\chi+\gamma\,\mathrm{Var}^{\mathrm{eff}}_t>\alpha_0\lambda/\Delta t$.
\end{theorem}

\begin{proposition}[Closed-loop stability in a homogeneous benchmark]
\label{prop:stability}
Assume the homogeneous closure \eqref{eq:mean_field_closure} holds and, for this benchmark calculation, that the mean belief tracks the fundamental ($m_t=v_t$). Then the mispricing $y_t:=p_t-v_t$ satisfies the linear SDE
\begin{equation}
\begin{aligned}
dy_t &= -\kappa_{\mathrm{eff}}\,y_t\,dt - \mu_v\,dt + \phi\,dZ_t - \sigma_v\,dW_t,\\
\kappa_{\mathrm{eff}} &:= \kappa\,\frac{\gamma\phi^2}{\gamma\phi^2-\lambda}.
\end{aligned}
\end{equation}
In particular, if $\gamma\phi^2>\lambda$ then $\kappa_{\mathrm{eff}}>0$ and $y_t$ is mean-reverting (Ornstein--Uhlenbeck). If additionally $\mu_v=0$, the stationary variance is $\Var(y_\infty)=(\phi^2+\sigma_v^2)/(2\kappa_{\mathrm{eff}})$.
\end{proposition}

\subsection{Uniqueness and Stability Conditions (Multiplicity and Blow-Ups)}

The following conditions rule out multiple equilibria and blow-ups. \textbf{Continuous-time LQG map} $\Phi$ (position-impact benchmark; Section~\ref{sec:lqg_constructive}): if $\gamma\phi^2>\lambda$ (positive denominator) and $q<1$, then $\Phi$ is a contraction on $(\mathcal X,\|\cdot\|_2)$, so the fixed point is unique and Picard iteration does not blow up (see Proposition~\ref{prop:lqg_constructive}). \textbf{Discrete-time per-step (baseline stock--flow):} if $0\le B_t<\Delta t$, the scalar stock--flow fixed point for $(\bar x_t,\bar u_t)$ is unique (Proposition~\ref{prop:step_fixed_point}); a sufficient uniform condition is \eqref{eq:global_contraction} (see Theorem~\ref{thm:contraction}). \textbf{Closed-loop paths (mispricing OU benchmark):} under the position-impact benchmark closure \eqref{eq:mean_field_closure}, if $\gamma\phi^2>\lambda$ then $\kappa_{\mathrm{eff}}>0$ and mispricing is mean-reverting (Proposition~\ref{prop:stability}). \textbf{Black-box MFG limit (benchmark):} under (H1)--(H5), the equilibrium flow is unique and $\bar\mu^N_t$ converges to $\mu_t$ at rate $O(1/N)$ (see Proposition~\ref{prop:mfg_limit_stability} in the mean-field limit subsection).

\subsection{Option B: Stationary Moments With Explicit \texorpdfstring{$k$}{k}-Dependence}

This subsection is a diagnostic benchmark built on the reduced-form position-impact closure \eqref{eq:mean_field_closure}; the baseline simulations use the discrete-time stock--flow equilibrium with uniqueness condition $0\le B_t<\Delta t$ (Proposition~\ref{prop:step_fixed_point}, Theorem~\ref{thm:contraction}).

Proposition~\ref{prop:stability} provides a transparent OU benchmark for mispricing, but it removes the overconfidence channel by imposing $m_t=v_t$. We now record complementary stationary-moment calculations in which $k$ enters explicitly through the steady-state filtering gain (Appendix~A).
Lemma~\ref{lem:riccati_convergence} shows that the perceived Riccati equation converges exponentially to the steady state, motivating the constant-gain approximation used below (interpretable as an asymptotic stationary regime after burn-in).

\begin{proposition}[Stationary disagreement and monotonicity in $k$ (constant-gain approximation)]
\label{prop:stationary_disagreement}
Assume a homogeneous overconfidence type $k$ and approximate the Kalman--Bucy gain by its steady-state value
$K^*(k)=\sigma_v/\sqrt{R'(k)}$, where $R'(k)=\sigma_c^2+\sigma_\epsilon^2/k$ (Appendix~A). Let
$m_t:=\int_0^1 \hat v_{i,t}\,di$ and define the cross-sectional deviation $\delta_{i,t}:=\hat v_{i,t}-m_t$.
Conditioning on the common filtration $\mathcal F_t^0$, the deviations satisfy the OU dynamics
\begin{equation}
d\delta_{i,t} = -K^*(k)\,\delta_{i,t}\,dt + K^*(k)\sigma_\epsilon\,dB_{i,t}.
\label{eq:delta_ou}
\end{equation}
Consequently, the conditional disagreement
$D_t:=\Var(\hat v_{i,t}\mid \mathcal F_t^0)=\Var(\delta_{i,t}\mid \mathcal F_t^0)$
admits the stationary value
\begin{equation}
D_\infty(k)=\frac{K^*(k)\sigma_\epsilon^2}{2}
=\frac{\sigma_v\sigma_\epsilon^2}{2\sqrt{\sigma_c^2+\sigma_\epsilon^2/k}},
\label{eq:D_infty}
\end{equation}
which is strictly increasing in $k$.
\end{proposition}

\begin{proposition}[Stationary mean-belief tracking error (constant-gain approximation)]
\label{prop:tracking_error}
Under the same constant-gain approximation, define the mean-belief tracking error $e_t:=m_t-v_t$. Then
\begin{equation}
de_t = -K^*(k)\,e_t\,dt + K^*(k)\sigma_c\,dU_t - \sigma_v\,dW_t,
\label{eq:e_ou}
\end{equation}
and the stationary variance is
\begin{equation}
\Var(e_\infty)=\frac{(K^*(k))^2\sigma_c^2+\sigma_v^2}{2K^*(k)}
=\frac{K^*(k)\sigma_c^2}{2}+\frac{\sigma_v^2}{2K^*(k)}.
\label{eq:Var_e_infty}
\end{equation}
In particular, $k$ enters the stationary tracking quality through $K^*(k)$ and $\Var(e_\infty)$ is strictly decreasing in $k$ (and converges to $\sigma_v\sigma_c$ as $k\to\infty$ when $\sigma_c>0$).
\end{proposition}

\begin{proposition}[Stationary mispricing with belief feedback (2D OU / Lyapunov form)]
\label{prop:2d_ou}
Assume $\mu_v=0$ and use the homogeneous closure \eqref{eq:mean_field_closure}. Under the constant-gain approximation above, the pair
$(y_t,e_t)$ with $y_t:=p_t-v_t$ satisfies the linear system
\begin{align}
dy_t &= -(\kappa+c)\,y_t\,dt + c\,e_t\,dt + \phi\,dZ_t - \sigma_v\,dW_t,\nonumber\\
de_t &= -K^*(k)\,e_t\,dt + K^*(k)\sigma_c\,dU_t - \sigma_v\,dW_t,
\label{eq:2d_ou}
\end{align}
where $c:=\lambda\kappa/(\gamma\phi^2-\lambda)$.
If $\gamma\phi^2>\lambda$ and $K^*(k)>0$, then the drift matrix
\[
A(k):=\begin{bmatrix}-(\kappa+c) & c\\ 0 & -K^*(k)\end{bmatrix}
\]
is Hurwitz and the stationary covariance $P(k)=\E[X_\infty X_\infty^\top]$ for $X_t=(y_t,e_t)$ is the unique solution of the Lyapunov equation
\begin{equation}
A(k)P(k)+P(k)A(k)^\top + Q(k) = 0,
\label{eq:lyapunov}
\end{equation}
where $Q(k)$ is the diffusion covariance implied by \eqref{eq:2d_ou}. Thus $\Var(y_\infty)$ depends on $k$ explicitly through $K^*(k)$.
\end{proposition}

\begin{corollary}[Closed-form stationary variances for Proposition~\ref{prop:2d_ou}]
\label{cor:2d_ou_closed_form}
In the setting of Proposition~\ref{prop:2d_ou}, write $K:=K^*(k)$ and define the noise covariances
$q_{11}:=\phi^2+\sigma_v^2$, $q_{12}:=\sigma_v^2$, and $q_{22}:=K^2\sigma_c^2+\sigma_v^2$.
Then the Lyapunov equation \eqref{eq:lyapunov} has the explicit solution $P=\begin{pmatrix}p_{11}&p_{12}\\p_{12}&p_{22}\end{pmatrix}$ with
\begin{equation}
\begin{aligned}
p_{22}&=\frac{q_{22}}{2K}=\frac{K\sigma_c^2}{2}+\frac{\sigma_v^2}{2K},\\
p_{12}&=\frac{c\,p_{22}+q_{12}}{\kappa+c+K},\\
p_{11}&=\frac{c\,p_{12}+q_{11}/2}{\kappa+c}.
\end{aligned}
\end{equation}
In particular, $\Var(y_\infty)=p_{11}$ depends on $k$ through $K^*(k)$.
\end{corollary}

\subsection{Connection to Numerical Results}

Section~VI simulates a large-$N$ particle system implementing \eqref{eq:anchored_impact}, Kalman belief updates (Appendix~A), and the feedback policy \eqref{eq:myopic_policy} (implemented in discrete time with an effective variance that includes belief uncertainty). Overconfidence enters only through belief updating (misperceived signal precision), and we report price/mispricing/volatility statistics across bull/bear drifts $\mu_v$ and overconfidence levels $k$.


%=================================================================
% Section V: Proofs
%=================================================================
\section{PROOFS}
\label{sec:proofs}

%------------------------------
\subsection*{Proof of Proposition~\ref{prop:static_comparative}}
\begin{proof}
From the static price equation \eqref{eq:static_price}, the coefficient on $v$ is
\[
a(k)=\frac{\lambda \tau'_{\epsilon,0}}{\gamma+\lambda(\tau_{v,0}+\tau'_{\epsilon,0})},
\qquad \tau'_{\epsilon,0}=\frac{k}{\sigma_{\epsilon,0}^2},\quad \tau_{v,0}=\frac{1}{\sigma_{v,0}^2}.
\]
Set $u:=\tau'_{\epsilon,0}>0$ and define $f(u):=\frac{\lambda u}{\gamma+\lambda(\tau_{v,0}+u)}$.
Then $a(k)=f(\tau'_{\epsilon,0})$ and
\[
f'(u)=\frac{\lambda(\gamma+\lambda\tau_{v,0})}{\big(\gamma+\lambda(\tau_{v,0}+u)\big)^2}>0,
\]
so $f$ is strictly increasing in $u$, hence $a(k)$ is strictly increasing in $k$.
Moreover, $a(k)>0$ is immediate since all parameters are positive, and
\[
a(k)<1 \iff \lambda u < \gamma+\lambda(\tau_{v,0}+u) \iff 0<\gamma+\lambda\tau_{v,0},
\]
which holds. Thus $0<a(k)<1$ and $a(k)$ increases with $k$.
\end{proof}

%------------------------------
\subsection*{Proof of Proposition~\ref{prop:myopic_cara}}
\begin{proof}
Fix a decision interval $\Delta t$ and condition on $\mathcal F_t^i$.
By the proposition's assumptions, $\Delta p_{t+\Delta t}$ is Gaussian with
\[
\mu_{i,t}:=\E_i[\Delta p_{t+\Delta t}]=\big(\kappa(\hat v_{i,t}-p_t)+\lambda\bar u_t\big)\Delta t,
\qquad
\Var_i(\Delta p_{t+\Delta t})=\sigma^2_{p,i,t}\Delta t.
\]
The one-step wealth increment is
\[
\Delta W_{i,t}=x_{i,t}\Delta p_{t+\Delta t}-\frac{\chi}{2}x_{i,t}^2\Delta t,
\]
hence $\Delta W_{i,t}$ is Gaussian conditional on $\mathcal F_t^i$ with
\[
\E_i[\Delta W_{i,t}]=x_{i,t}\mu_{i,t}-\frac{\chi}{2}x_{i,t}^2\Delta t,
\qquad
\Var_i(\Delta W_{i,t})=x_{i,t}^2\sigma^2_{p,i,t}\Delta t.
\]
For CARA utility $U(w)=-e^{-\gamma w}$ and Gaussian $X\sim N(m,s^2)$, the standard identity gives
$\E[-e^{-\gamma X}]=-\exp(-\gamma m+(\gamma^2/2)s^2)$.
Applying this with $X=W_{i,t}+\Delta W_{i,t}$ (where $W_{i,t}$ is $\mathcal F_t^i$-measurable),
maximizing conditional expected utility over $x_{i,t}$ is equivalent to maximizing the certainty-equivalent quadratic
\[
x_{i,t}\mu_{i,t}-\frac{\gamma}{2}x_{i,t}^2\sigma^2_{p,i,t}\Delta t-\frac{\chi}{2}x_{i,t}^2\Delta t.
\]
The first-order condition yields
$\mu_{i,t}-(\gamma\sigma^2_{p,i,t}+\chi)x_{i,t}\Delta t=0$, so
\[
x^\star_{i,t}=\frac{\mu_{i,t}/\Delta t}{\chi+\gamma\sigma^2_{p,i,t}}
=\frac{\kappa(\hat v_{i,t}-p_t)+\lambda\bar u_t}{\chi+\gamma\sigma^2_{p,i,t}},
\]
which is \eqref{eq:myopic_policy}.
\end{proof}

%------------------------------
\subsection*{Proof of Proposition~\ref{prop:step_fixed_point}}
\begin{proof}
Work on the discrete grid with step $\Delta t$ and the Euler increment
\[
\begin{aligned}
\Delta p_{t+1}:=p_{t+1}-p_t
&=\kappa(v_t-p_t)\Delta t+\lambda \bar u_t\Delta t\\
&\quad +\phi\sqrt{\Delta t}\,\varepsilon^\eta_{t+1},\\
\varepsilon^\eta_{t+1}&\sim N(0,1).
\end{aligned}
\]
Conditional on investor $i$'s information $\mathcal F_t^i$, we have $\mathbb E[v_t\mid\mathcal F_t^i]=\hat v_{i,t}$ and
$\mathrm{Var}(v_t\mid\mathcal F_t^i)=\Sigma_{i,t}$. Hence
\[
\mathbb E[\Delta p_{t+1}\mid\mathcal F_t^i]
=\big(\kappa(\hat v_{i,t}-p_t)+\lambda \bar u_t\big)\Delta t,
\]
and, using independence of $\varepsilon^\eta_{t+1}$ from $\mathcal F_t^i$ and from $v_t$,
\[
\mathrm{Var}(\Delta p_{t+1}\mid\mathcal F_t^i)
\begin{aligned}
&=\kappa^2\,\mathrm{Var}(v_t\mid\mathcal F_t^i)(\Delta t)^2+\phi^2\Delta t\\
&=\kappa^2\Sigma_{i,t}(\Delta t)^2+\phi^2\Delta t.
\end{aligned}
\]
Under the one-step myopic CARA approximation (Proposition~\ref{prop:myopic_cara}), the investor chooses $x_{i,t}$ to maximize the certainty-equivalent quadratic
\[
 x_{i,t}\,\mathbb E[\Delta p_{t+1}\mid\mathcal F_t^i]
 -\frac{\gamma}{2}x_{i,t}^2\,\mathrm{Var}(\Delta p_{t+1}\mid\mathcal F_t^i)
 -\frac{\chi}{2}x_{i,t}^2\,\Delta t.
\]
Substituting the conditional moments above and dividing by the common factor $\Delta t>0$ gives the equivalent problem
\[
\max_{x_{i,t}}\;
 x_{i,t}\big(\kappa(\hat v_{i,t}-p_t)+\lambda \bar u_t\big)
 -\frac12\Big(\gamma\big(\phi^2+\kappa^2\Sigma_{i,t}\Delta t\big)+\chi\Big)x_{i,t}^2.
\]
The unique maximizer is
\[
 x_{i,t}^*=\frac{\kappa(\hat v_{i,t}-p_t)+\lambda \bar u_t}{\chi+\gamma(\phi^2+\kappa^2\Sigma_{i,t}\Delta t)}.
\]
Including the optional scale $\alpha_0>0$ (risk-budget/position-size parameter) yields \eqref{eq:impl_policy} with
\[
 \mathrm{Var}^{\mathrm{eff}}_{i,t}:=\phi^2+\kappa^2\Sigma_{i,t}\Delta t,
 \]
 
Write \eqref{eq:impl_policy} as $x_{i,t}=a_{i,t}+b_{i,t}\bar u_t$ with
\[
 a_{i,t}:=\frac{\alpha_0\kappa(\hat v_{i,t}-p_t)}{\chi+\gamma \mathrm{Var}^{\mathrm{eff}}_{i,t}},
 \qquad
 b_{i,t}:=\frac{\alpha_0\lambda}{\chi+\gamma \mathrm{Var}^{\mathrm{eff}}_{i,t}}.
\]
Averaging over $i=1,\dots,N$ gives $\bar x_t=A_t+B_t\bar u_t$ where $A_t,B_t$ are as defined in the paper.
Using $\bar u_t=(\bar x_t-\bar x_{t^-})/\Delta t$ gives
\[
\bar x_t = A_t+\frac{B_t}{\Delta t}\bar x_t-\frac{B_t}{\Delta t}\bar x_{t^-}.
\]
Rearranging yields $\big(1-B_t/\Delta t\big)\bar x_t=A_t-(B_t/\Delta t)\bar x_{t^-}$.
If $0\le B_t<\Delta t$, the unique solution is
\[
\bar x_t=\frac{A_t-(B_t/\Delta t)\bar x_{t^-}}{1-B_t/\Delta t},
\qquad
\bar u_t=\frac{\bar x_t-\bar x_{t^-}}{\Delta t},
\]
as claimed.
\end{proof}

%------------------------------
\subsection*{Proof of Theorem~\ref{thm:contraction}}
\begin{proof}
By Proposition~\ref{prop:step_fixed_point}, for fixed $(\hat v_{i,t},\Sigma_{i,t})_{i=1}^N$, $p_t$, and $\bar x_{t^-}$, the per-step map is affine:
\[
 \Phi_t(\bar x)=A_t+\frac{B_t}{\Delta t}\bar x-\frac{B_t}{\Delta t}\bar x_{t^-}.
\]
For any $\bar x,\bar x'$,
\[
 |\Phi_t(\bar x)-\Phi_t(\bar x')|=\frac{|B_t|}{\Delta t}\,|\bar x-\bar x'|.
\]
Thus if $0\le B_t<\Delta t$, $\Phi_t$ is a contraction on $\mathbb R$ with contraction constant $B_t/\Delta t$ and has a unique fixed point,
given in closed form in Proposition~\ref{prop:step_fixed_point}.

For the sufficient uniform condition: since $\Sigma_{i,t}\ge 0$ and $\Delta t>0$, we have
$\mathrm{Var}^{\mathrm{eff}}_{i,t}=\phi^2+\kappa^2\Sigma_{i,t}\Delta t\ge \phi^2$, hence
$D_{i,t}=\chi+\gamma\mathrm{Var}^{\mathrm{eff}}_{i,t}\ge \chi+\gamma\phi^2$ and therefore
\[
 B_t=\frac{1}{N}\sum_{i=1}^N \frac{\alpha_0\lambda}{D_{i,t}}
 \le \frac{1}{N}\sum_{i=1}^N \frac{\alpha_0\lambda}{\chi+\gamma \phi^2}
 =\frac{\alpha_0\lambda}{\chi+\gamma\phi^2}.
\]
Therefore $\alpha_0\lambda/(\Delta t(\chi+\gamma\phi^2))<1$ implies $B_t<\Delta t$ for all $t$.
\end{proof}

%------------------------------
\subsection*{Proof of Theorem~\ref{thm:well_posed} (Well-posedness of the discrete-time particle system)}
\begin{proof}
Fix $T\in\mathbb N$, $\Delta t>0$, and fix a realization of all common and idiosyncratic shock sequences used by the algorithm.
We show by induction on $t=0,1,\dots,T$ that the state
\[
S_t=(v_t,p_t,\{(\hat v_{i,t},\Sigma_{i,t})\}_{i=1}^N,\bar x_t)
\]
is uniquely determined and adapted.

At $t=0$, $(v_0,p_0,\{\hat v_{i,0},\Sigma_{i,0}\}_{i=1}^N)$ are given by assumption (A1). Suppose the system is uniquely
defined up to time $t$. Then:
\begin{enumerate}[leftmargin=2em]
\item By (A2), the discrete-time Kalman recursion in Appendix A defines each $(\hat v_{i,t},\Sigma_{i,t})$ uniquely from
$(\hat v_{i,t-1},\Sigma_{i,t-1})$ and the realized private signal at time $t$ (which is a measurable function of the fixed
shock sequences and $v_t$).
\item Given $(\hat v_{i,t},\Sigma_{i,t})_{i=1}^N$ and $p_t$, the per-step demand map $\Phi_t$ is well-defined by Proposition~\ref{prop:step_fixed_point}.
By (A3) (e.g.\ $B_t<\Delta t$), $\Phi_t$ has a unique fixed point $\bar x_t$ (and hence $\bar u_t$), and then each $x_{i,t}$ is uniquely determined.
\item Finally, by (A4), the Euler updates for $v_{t+1}$ and $p_{t+1}$ are explicit measurable functions of $(v_t,p_t)$, $\bar u_t$,
and the realized shocks at time $t+1$, hence define unique $(v_{t+1},p_{t+1})$.
\end{enumerate}
Thus $S_{t+1}$ is uniquely determined from $S_t$ and the fixed shocks, and is adapted to the natural filtration generated by
the shocks. Induction completes the proof, yielding a unique adapted solution path on $\{0,1,\dots,T\}$.
\end{proof}

%------------------------------
\subsection*{Proof of Proposition~\ref{prop:stationary_disagreement} (Constant-gain approximation)}
\begin{proof}
Assume a homogeneous type $k$ and adopt the constant-gain approximation $K_t\equiv K^*(k)=:K>0$.
From Appendix A, the (subjective) Kalman--Bucy filter is
\[
d\hat v_{i,t}=\mu_v\,dt+K\big(d\xi_{i,t}-\hat v_{i,t}\,dt\big).
\]
Using the observation equation \eqref{eq:xi_dyn_ct},
\[
d\xi_{i,t}=v_t\,dt+\sigma_c\,dU_t+\sigma_\epsilon\,dB_{i,t},
\]
so substituting gives
\[
\begin{aligned}
d\hat v_{i,t}
&=\mu_v\,dt+K\big((v_t-\hat v_{i,t})\,dt+\sigma_c\,dU_t+\sigma_\epsilon\,dB_{i,t}\big)\\
&= \big(\mu_v+K(v_t-\hat v_{i,t})\big)\,dt+K\sigma_c\,dU_t+K\sigma_\epsilon\,dB_{i,t}.
\end{aligned}
\]
Let $m_t:=\int_0^1 \hat v_{i,t}\,di$ be the cross-sectional mean belief. Under a continuum (or large-$N$) approximation,
the idiosyncratic noises average out, i.e.\ $\int_0^1 dB_{i,t}\,di=0$ in the sense of laws/LLN, while the common term
$U_t$ persists. Averaging the SDE over $i$ yields
\[
dm_t=\big(\mu_v+K(v_t-m_t)\big)\,dt+K\sigma_c\,dU_t.
\]
Define the deviation $\delta_{i,t}:=\hat v_{i,t}-m_t$. Subtracting the $m_t$ equation from the $\hat v_{i,t}$ equation cancels
the common-noise terms and gives
\[
\begin{aligned}
d\delta_{i,t}
&= -K(\hat v_{i,t}-m_t)\,dt + K\sigma_\epsilon\,dB_{i,t}\\
&= -K\delta_{i,t}\,dt+K\sigma_\epsilon\,dB_{i,t},
\end{aligned}
\]
which is \eqref{eq:delta_ou}.

Now condition on the common filtration $\mathcal F_t^0:=\sigma(W_s,U_s,Z_s:s\le t)$. Conditional on $\mathcal F_t^0$,
$\delta_{i,t}$ solves an OU SDE driven only by the idiosyncratic Brownian motion $B_{i,t}$, so its conditional variance
$D_t:=\mathrm{Var}(\delta_{i,t}\mid \mathcal F_t^0)=\mathbb E[\delta_{i,t}^2\mid \mathcal F_t^0]$ satisfies (It\^o + tower property)
\[
\begin{aligned}
d(\delta_{i,t}^2)
&=2\delta_{i,t}\,d\delta_{i,t}+(d\delta_{i,t})^2\\
&=\big(-2K\delta_{i,t}^2+K^2\sigma_\epsilon^2\big)\,dt+2K\sigma_\epsilon\delta_{i,t}\,dB_{i,t}.
\end{aligned}
\]
Taking $\mathcal F_t^0$-conditional expectation kills the martingale term and yields the linear ODE
\[
\frac{d}{dt}D_t = -2K D_t + K^2\sigma_\epsilon^2.
\]
Hence $D_t\to D_\infty:=\frac{K^2\sigma_\epsilon^2}{2K}=\frac{K\sigma_\epsilon^2}{2}$, giving \eqref{eq:D_infty}.

Finally, with $R'(k)=\sigma_c^2+\sigma_\epsilon^2/k$ we have $K^*(k)=\sigma_v/\sqrt{R'(k)}$, which is strictly increasing in $k$
because $R'(k)$ is strictly decreasing in $k$; therefore $D_\infty(k)=(K^*(k)\sigma_\epsilon^2)/2$ is strictly increasing in $k$.
\end{proof}

%------------------------------
\subsection*{Proof of Proposition~\ref{prop:tracking_error} (Constant-gain approximation)}
\begin{proof}
Under the constant-gain approximation (same setting as Proposition~\ref{prop:stationary_disagreement}), we derived
\[
dm_t=\big(\mu_v+K(v_t-m_t)\big)\,dt+K\sigma_c\,dU_t.
\]
The fundamental satisfies $dv_t=\mu_v\,dt+\sigma_v\,dW_t$. Define $e_t:=m_t-v_t$.
Then
\[
\begin{aligned}
de_t
&=dm_t-dv_t\\
&=\big(\mu_v+K(v_t-m_t)\big)\,dt+K\sigma_c\,dU_t-\mu_v\,dt-\sigma_v\,dW_t\\
&=-K e_t\,dt+K\sigma_c\,dU_t-\sigma_v\,dW_t,
\end{aligned}
\]
which is \eqref{eq:e_ou}.

Since $U$ and $W$ are independent Brownian motions, $e_t$ is an OU process with mean-reversion rate $K$ and diffusion variance rate
$K^2\sigma_c^2+\sigma_v^2$. The stationary variance is the standard OU value
\[
\mathrm{Var}(e_\infty)=\frac{K^2\sigma_c^2+\sigma_v^2}{2K}
=\frac{K\sigma_c^2}{2}+\frac{\sigma_v^2}{2K},
\]
which is \eqref{eq:Var_e_infty}.

To see monotonicity in $k$, view $\Var(e_\infty)$ as a function of $K>0$:
\[
g(K):=\frac{K\sigma_c^2}{2}+\frac{\sigma_v^2}{2K},
\qquad
g'(K)=\frac{\sigma_c^2}{2}-\frac{\sigma_v^2}{2K^2}.
\]
If $\sigma_c=0$, then $g(K)=\sigma_v^2/(2K)$ is strictly decreasing in $K$, and since $K^*(k)$ is strictly increasing in $k$,
$\Var(e_\infty)$ is strictly decreasing in $k$.
If $\sigma_c>0$, then $g'(K)<0$ on $K\in(0,\sigma_v/\sigma_c)$ and $K^*(k)=\sigma_v/\sqrt{R'(k)}\le \sigma_v/\sigma_c$ since
$R'(k)=\sigma_c^2+\sigma_\epsilon^2/k\ge \sigma_c^2$. Thus $g$ is decreasing along the range of $K^*(k)$, so $\Var(e_\infty)$ is
strictly decreasing in $k$.
\end{proof}

%------------------------------
\subsection*{Proof of Proposition~\ref{prop:2d_ou} (2D OU / Lyapunov form)}
\begin{proof}
Assume $\mu_v=0$ and the homogeneous closure \eqref{eq:mean_field_closure}:
\[
\bar x_t=\frac{\kappa}{\gamma\phi^2-\lambda}(m_t-p_t).
\]
Write $y_t:=p_t-v_t$ and $e_t:=m_t-v_t$. Using $m_t-p_t=(m_t-v_t)-(p_t-v_t)=e_t-y_t$, the price SDE \eqref{eq:anchored_impact} becomes
\[
\begin{aligned}
dp_t
&=\kappa(v_t-p_t)\,dt+\lambda \bar x_t\,dt+\phi\,dZ_t\\
&= -\kappa y_t\,dt + \frac{\lambda\kappa}{\gamma\phi^2-\lambda}(e_t-y_t)\,dt+\phi\,dZ_t.
\end{aligned}
\]
Let $c:=\lambda\kappa/(\gamma\phi^2-\lambda)$. Then
\[
dp_t=-(\kappa+c)y_t\,dt+c e_t\,dt+\phi\,dZ_t.
\]
Since $\mu_v=0$, $dv_t=\sigma_v\,dW_t$, hence
\[
dy_t=dp_t-dv_t=-(\kappa+c)y_t\,dt+c e_t\,dt+\phi\,dZ_t-\sigma_v\,dW_t.
\]
From Proposition~\ref{prop:tracking_error}, the mean-belief tracking error satisfies
\[
de_t=-K e_t\,dt+K\sigma_c\,dU_t-\sigma_v\,dW_t,
\qquad K:=K^*(k)>0.
\]
Stacking $X_t=(y_t,e_t)^\top$ yields the linear system \eqref{eq:2d_ou} with drift matrix
\[
A(k)=\begin{pmatrix}-(\kappa+c) & c\\ 0 & -K\end{pmatrix}
\]
and diffusion matrix $G(k)$ as stated.

If $\gamma\phi^2>\lambda$ then $c>0$ and $\kappa+c>0$; also $K>0$ by definition. The eigenvalues of $A(k)$ are
$-(\kappa+c)$ and $-K$, hence strictly negative: $A(k)$ is Hurwitz. For a Hurwitz $A$, the continuous-time Lyapunov equation
\[
A P + P A^\top + G G^\top = 0
\]
has a unique symmetric positive semidefinite solution $P$, which equals the stationary covariance $P=\mathbb E[X_\infty X_\infty^\top]$.
This gives \eqref{eq:lyapunov} and shows $\mathrm{Var}(y_\infty)$ depends on $k$ through $K^*(k)$.
\end{proof}

%------------------------------
\subsection*{Proof of Corollary~\ref{cor:2d_ou_closed_form} (Closed-form stationary variances)}
\begin{proof}
Let $A=\begin{pmatrix}-(\kappa+c) & c\\ 0 & -K\end{pmatrix}$ and $Q:=GG^\top$.
Because $dZ_t,dU_t,dW_t$ are independent, from \eqref{eq:2d_ou} we have
\[
Q=\begin{pmatrix}
\phi^2+\sigma_v^2 & \sigma_v^2\\
\sigma_v^2 & K^2\sigma_c^2+\sigma_v^2
\end{pmatrix}
 =: \begin{pmatrix} q_{11} & q_{12}\\ q_{12} & q_{22}\end{pmatrix}.
\]
Write $P=\begin{pmatrix}p_{11}&p_{12}\\ p_{12}&p_{22}\end{pmatrix}$.
Compute
\[
\begin{aligned}
AP&=\begin{pmatrix}
-(\kappa+c)p_{11}+c p_{12} & -(\kappa+c)p_{12}+c p_{22}\\
- K p_{12} & -K p_{22}
\end{pmatrix},\\[2pt]
PA^\top&=\begin{pmatrix}
-(\kappa+c)p_{11}+c p_{12} & -K p_{12}\\
-(\kappa+c)p_{12}+c p_{22} & -K p_{22}
\end{pmatrix}.
\end{aligned}
\]
Thus the Lyapunov equation $AP+PA^\top+Q=0$ yields the entrywise system:
\begin{align*}
(2,2):\;& -2K p_{22}+q_{22}=0,\\
(1,2):\;& -(\kappa+c+K)p_{12}+c p_{22}+q_{12}=0,\\
(1,1):\;& -2(\kappa+c)p_{11}+2c p_{12}+q_{11}=0.
\end{align*}
Solving sequentially gives
\[
p_{22}=\frac{q_{22}}{2K}=\frac{K^2\sigma_c^2+\sigma_v^2}{2K}
=\frac{K\sigma_c^2}{2}+\frac{\sigma_v^2}{2K},
\]
\[
p_{12}=\frac{c p_{22}+q_{12}}{\kappa+c+K},
\qquad
p_{11}=\frac{c p_{12}+q_{11}/2}{\kappa+c},
\]
which matches the closed-form expressions in Corollary~\ref{cor:2d_ou_closed_form}.
\end{proof}

%------------------------------
\subsection*{Proof of Proposition~\ref{prop:stability} (Homogeneous closed-loop stability benchmark)}
\begin{proof}
Assume the homogeneous closure \eqref{eq:mean_field_closure} holds and, for this benchmark, $m_t=v_t$. Then
\[
\bar x_t=\frac{\kappa}{\gamma\phi^2-\lambda}(v_t-p_t),
\]
and substituting into \eqref{eq:anchored_impact} gives
\[
\begin{aligned}
dp_t
&=\kappa(v_t-p_t)\,dt+\lambda\frac{\kappa}{\gamma\phi^2-\lambda}(v_t-p_t)\,dt+\phi\,dZ_t\\
&= -\kappa_{\mathrm{eff}}(p_t-v_t)\,dt+\phi\,dZ_t,
\end{aligned}
\]
where
\[
\kappa_{\mathrm{eff}}:=\kappa\Big(1+\frac{\lambda}{\gamma\phi^2-\lambda}\Big)
=\kappa\frac{\gamma\phi^2}{\gamma\phi^2-\lambda}.
\]
Define $y_t:=p_t-v_t$. Using $dv_t=\mu_v\,dt+\sigma_v\,dW_t$,
\[
dy_t=dp_t-dv_t=-\kappa_{\mathrm{eff}}y_t\,dt-\mu_v\,dt+\phi\,dZ_t-\sigma_v\,dW_t,
\]
as claimed. If $\gamma\phi^2>\lambda$ then $\kappa_{\mathrm{eff}}>0$ and $y_t$ is an OU process.
When $\mu_v=0$, $y_t$ has diffusion variance rate $\phi^2+\sigma_v^2$ and mean-reversion rate $\kappa_{\mathrm{eff}}$,
so its stationary variance is
\[
\mathrm{Var}(y_\infty)=\frac{\phi^2+\sigma_v^2}{2\kappa_{\mathrm{eff}}}.
\]
\end{proof}


\section{Numerical Examples (technical extract)}
\label{sec:numerical}
% Review version (technical extract) of Section 6: Numerical Examples.
% Keeps only the algorithmic/procedural content needed to parse the model's implementation bridge,
% plus label anchors referenced from the unchanged introduction.

\paragraph{Scheme (map form).}
Let $S_t=(v_t,p_t,\{(\hat v_{i,t},\Sigma_{i,t},x_{i,t^-})\}_{i=1}^N)$. For fixed shocks (common and idiosyncratic) the particle system defines a measurable one-step map
\begin{equation}
S_{t+1}=\Psi_{\Delta t,N}(S_t;\text{shocks}),
\end{equation}
obtained by composing: (a) the Kalman update (Appendix~A), (b) the per-step mean-field fixed point (Theorem~\ref{thm:contraction}), and (c) the Euler updates for $(v,p)$. Under the contraction condition \eqref{eq:global_contraction}, the fixed-point step is stable and the full map is well-defined (Theorem~\ref{thm:well_posed}).

\begin{algorithm}[t]
\caption{Particle approximation (one scenario)}
\label{alg:particle}
\begin{algorithmic}[1]
\STATE Fix parameters, scenario $(\mu_v,k)$, horizon $T$, population size $N$, and seeds $m=1,\dots,M$.
\FOR{$m=1$ to $M$}
  \STATE Simulate one realization of common shocks $(\Delta W_t,\Delta U_t,\Delta Z_t)_{t=0}^{T-1}$.
  \STATE Initialize $(v_0,p_0)$ and agent beliefs $(\hat v_{i,0},\Sigma_{i,0})$ for $i=1,\dots,N$.
  \STATE Initialize inventories $x_{i,0^-}$ and wealth $W_{i,0}$ for $i=1,\dots,N$.
  \FOR{$t=0$ to $T-1$}
    \STATE For each agent $i$, draw idiosyncratic signal noise $\Delta B_{i,t}$ and update $(\hat v_{i,t},\Sigma_{i,t})$ via the (perceived) Kalman recursion.
    \STATE Compute $(\bar x_t,\bar u_t)$ from the per-step stock--flow fixed point using pre-trade inventories (Proposition~\ref{prop:step_fixed_point}).
    \STATE For each agent $i$, set $x_{i,t}=x^\star_{i,t}$ from the myopic rule given $\bar u_t$ and compute trading rate $u_{i,t}=(x_{i,t}-x_{i,t^-})/\Delta t$.
    \STATE Update $(v_{t+1},p_{t+1})$ using the Euler discretization of \eqref{eq:anchored_impact} (flow-based impact).
    \STATE Update wealth $W_{i,t+1}=W_{i,t}+x_{i,t}\Delta p_{t+1}-(\chi/2)x_{i,t}^2\Delta t$ and set $x_{i,t+1^-}=x_{i,t}$.
  \ENDFOR
  \STATE Record seed-level time series (price, fundamental, returns, mispricing, demand moments).
\ENDFOR
\STATE Report scenario figures/tables by averaging statistics across seeds (Monte Carlo over the common and idiosyncratic shocks).
\end{algorithmic}
\end{algorithm}

% Parameter table for model and numerical settings (particle simulations)

\begin{table}[t]
\centering
\scriptsize
\setlength{\tabcolsep}{3pt}
\caption{Model and numerical parameters.}
\label{tab:params}
\begin{tabular}{lp{1.8cm}cc}
\toprule
\textbf{Parameter} & \textbf{Description} & \textbf{Baseline Value} & \textbf{Units} \\
\midrule
\multicolumn{4}{l}{\textit{Static Benchmark (Sec.~2.2)}} \\
\midrule
$\mu_{v,0}$ & Prior mean of $v$ & --- & value \\
$\sigma_{v,0}$ & Prior sd of $v$ & --- & value \\
$\sigma_{\epsilon,0}$ & Signal-noise sd & --- & value \\
\midrule
\multicolumn{4}{l}{\textit{Dynamic Model (Sec.~2.3+)}} \\
\midrule
$\gamma$ & Risk aversion & $1.0$ & --- \\
$\chi$ & Inventory holding cost & $0.05$ & value/(share$^2\cdot$step) \\
$k$ & Overconfi\-dence (precision scale) & $1.0$, $3.0$ & --- \\
$\mu_v$ & Fundamental drift (bull/bear) & $\pm 0.02$ & value/step \\
$\sigma_v$ & Fundamental diffusion vol. & $0.10$ & value/$\sqrt{\text{step}}$ \\
$\sigma_c$ & Common signal component & $0.10$ & value/$\sqrt{\text{step}}$ \\
$\sigma_\epsilon$ & Idiosyncratic signal loading & $0.80$ & value/$\sqrt{\text{step}}$ \\
$\kappa$ & Fundamental anchoring & $0.005$ & 1/step \\
$\lambda$ & Impact strength & $0.20$ & value/share \\
$\phi$ & Price/noise-trading volatility & $0.50$ & value/$\sqrt{\text{step}}$ \\
\midrule
\multicolumn{4}{l}{\textit{Numerical Parameters}} \\
\midrule
$N$ & Number of agents & $200$ & --- \\
$T$ & Horizon (steps) & $400$ & step \\
$M$ & Random seeds (replications) & $50$ & --- \\
$p_0$ & Initial price & $0.0$ & value \\
$v_0$ & Initial fundamental & $0.0$ & value \\
\bottomrule
\end{tabular}
\end{table}


\subsection{Parameter identification and calibration strategy}
\label{sec:param_identification}

While full empirical estimation is beyond the scope of this paper, we provide an identification argument showing which observable moments pin down the key structural parameters $(\kappa,\lambda,\phi,\sigma_\epsilon,\sigma_v)$ in principle. This clarifies what data would be needed for calibration and establishes the model's testability.

\paragraph{Identification via observable moments.}
Under the myopic mean-field equilibrium, the following moments are observable (or proxied) from price and volume data:

\begin{enumerate}[label=(\arabic*)]
\item \textbf{Return volatility $\mathrm{std}(r_t)$:} In the model, ``return'' is the price increment $r_t := \Delta p_t$ (value units). Empirically one often works with log returns $\Delta\log P_t$; for small moves, $\Delta\log P_t \approx \Delta P_t/P_t$, so a simple scaling can map model increments to empirical log-return units when needed. From \eqref{eq:anchored_impact}, the price increment variance is
\[
\Var(\Delta p_{t+1}\mid \mathcal F_t) \approx
\Big(\kappa^2\,\Var(v_t-p_t\mid \mathcal F_t)+\lambda^2\,\Var(\bar x_t\mid \mathcal F_t)+\phi^2\Big)\Delta t
\;+\;\text{cross terms},
\]
In steady state, $\mathrm{std}(r_t) = \mathrm{std}(\Delta p_t)$ depends on $\phi$ (direct noise trading), $\lambda$ (impact amplification), and $\kappa$ (mean reversion strength). The ratio $\lambda/\kappa$ can be identified from the relative contribution of demand-driven vs.\ fundamental-driven price variance.

\item \textbf{Mean-reversion speed:} The autocorrelation of price deviations from a moving average of fundamentals (or a proxy like a slow-moving average of prices) identifies $\kappa$. Specifically, under \eqref{eq:anchored_impact}, the half-life of price deviations is approximately $(\ln 2)/\kappa$, so $\kappa$ can be estimated from the decay rate of price-f fundamental autocorrelations.
\begin{equation}
h = \frac{\ln 2}{\kappa},\qquad \kappa = \frac{\ln 2}{h}.
\label{eq:half_life_conversion}
\end{equation}

\item \textbf{Impact coefficient $\lambda$:} The contemporaneous correlation between order flow (aggregate demand proxy) and price changes, normalized by the variance of order flow, identifies $\lambda$ up to the mean-field scaling. Volume proxies (e.g., signed volume or turnover) can serve as demand proxies.

\item \textbf{Fundamental volatility $\sigma_v$:} The long-run variance of fundamentals (estimated from a slow-moving filter of prices or from fundamental proxies like earnings/dividends) identifies $\sigma_v$. Alternatively, $\sigma_v$ can be identified from the steady-state variance of the Kalman filter innovation process.

\item \textbf{Signal noise $\sigma_\epsilon$:} The cross-sectional dispersion in beliefs (proxied by analyst forecast dispersion or survey-based disagreement) relative to the fundamental variance identifies the signal-to-noise ratio $\sigma_v/\sigma_\epsilon$. Given $\sigma_v$, this pins down $\sigma_\epsilon$.
\end{enumerate}

\paragraph{Identification chain (estimated vs.\ inferred vs.\ normalized).}
From price-only data, one can (a) estimate $\kappa$ from a half-life diagnostic (item 2) and (b) calibrate the effective price/noise-trading volatility $\phi$ to match return volatility (item 1, after the unit mapping to log-return units if desired).
Identifying $\lambda$ itself requires an order-flow proxy (item 3); given $\lambda$, the corresponding order-flow noise scale is inferred mechanically as $\sigma_\eta=\phi/\lambda$ (Section~II, microstructure sketch).
In our numerical work, $\lambda$ is treated as a normalization choice (Table~\ref{tab:params}), and we therefore interpret $\phi$ as the empirically pinned-down price-noise scale conditional on that normalization.

\subsection{Parameter regimes and stress diagnostics}
\label{sec:meaningful_regions}

% Label-only anchors for tables referenced elsewhere; the corresponding numerical tables are omitted in this review build.
\refstepcounter{table}
\label{tab:price_filter}
\refstepcounter{table}
\label{tab:mispricing_stress}

\subsubsection{Targeted joint-move grid}
\label{subsec:joint_grid}

\subsection{Extensions and validation checks}
\label{subsec:extensions}


% Anchor a label for the (omitted) conclusion section so cross-references in the unchanged introduction compile.
\refstepcounter{section}
\label{sec:conclusion}

\appendix
\input{sections/appx_review}

\bibliographystyle{plainnat}
\bibliography{refs}

\end{document}
